\documentclass[11pt]{amsart}
\IfFileExists{lmodern.sty}%
  {\usepackage[T1]{fontenc}\usepackage{lmodern}\usepackage{microtype}}%
  {\usepackage[expansion=false]{microtype}}
\usepackage{amssymb,mathtools}
\usepackage{booktabs,array}
\usepackage{enumitem}
\emergencystretch=2.5em
\usepackage[colorlinks=true,linkcolor=blue,citecolor=blue,urlcolor=blue]{hyperref}

\newcommand{\FormalizationURL}{https://github.com/slitsey/catalan-irrationality-measure}
\newcommand{\FormalizationLink}{\url{\FormalizationURL}}

% ---------------------------------------------------------------------------
% Numerical constants. Each is the decimal truncation (or outward rounding) of a rational
% number of the Lean proof; the Lean names are in the comments.
\newcommand{\Bx}{409.99999}                             % 53 rho'/(eta - 212 delta)
\newcommand{\Rho}{296.586094}                           % rho' = K590.rhoEnd (879663/10000)
\newcommand{\RhoOld}{301.54258}                         % 58 c_2 + w (used only in the remark on rho)
\newcommand{\Ctwo}{3.596893}                            % secondCost'
\newcommand{\Cone}{3.549542}                            % firstCost'
\newcommand{\PsiZero}{87.9663}                          % Psi_0, weight plus binomial per column
\newcommand{\Utarget}{2.255858836}                      % K590.target
\newcommand{\Gain}{8.29824}                             % 58^2 (target0 - target)
\newcommand{\Uzero}{2.2583256}                          % target0
\newcommand{\NormCap}{0.7736205195}                     % normCap
\newcommand{\SupX}{-0.96450499}
\newcommand{\SupY}{-1.57499400}
\newcommand{\FieldCap}{-2.53949899}                     % fieldCap
\newcommand{\ResCap}{-0.0631062}                        % residualCap
\newcommand{\EtaInf}{7.98775}                           % real place (old) minus denominators, no gain
\newcommand{\EtaTot}{38.33924}                          % eta of the criterion
\newcommand{\Kappa}{1.382397}                           % 410/rho'
\newcommand{\Kcoef}{160}
\newcommand{\Ngeo}{75}
\newcommand{\Ncomplex}{41}
\newcommand{\Nblocks}{47, 38 and 23}
\newcommand{\Ncells}{373, 302 and 176}
\newcommand{\Mmargin}{2.986\cdot10^{-6}}
\newcommand{\TDelta}{6431.1338}                         % 58^2 (target0 - log 2 / 2)
\newcommand{\TDeltaOld}{6422.8355}                      % 58^2 (target - log 2 / 2)
\newcommand{\EtaExact}{38.34580}                        % TDelta minus the denominator exponent
\newcommand{\PayL}{53.0000849}                          % d + 4 d delta kappa
\newcommand{\PayR}{53.0000857}                          % eta kappa
\newcommand{\Digest}{ff307ca0287dc5224e268202cb0b09776e933652acd2f2f879492af28d3e381d}
\newcommand{\Chi}{22.05972}                             % area of the rank-deficiency profile
\newcommand{\ChiLean}{22.05325}                         % gain532_ge
\newcommand{\Aden}{6392.7880}                           % exponent of the denominators
\newcommand{\Wcap}{52.788}
\newcommand{\NcellsSigned}{\textbf{??}}                 % filled in after the cell build
% ---------------------------------------------------------------------------

\newtheorem{theorem}{Theorem}[section]
\newtheorem{lemma}[theorem]{Lemma}
\newtheorem{proposition}[theorem]{Proposition}
\newtheorem{corollary}[theorem]{Corollary}
\theoremstyle{definition}
\newtheorem{definition}[theorem]{Definition}
\theoremstyle{remark}
\newtheorem{remark}[theorem]{Remark}

\numberwithin{equation}{section}
\setcounter{tocdepth}{1}

\newcommand{\Z}{\mathbb Z}
\newcommand{\Q}{\mathbb Q}
\newcommand{\R}{\mathbb R}
\newcommand{\N}{\mathbb N}
\newcommand{\C}{\mathbb C}
\newcommand{\F}{\mathbb F}
\newcommand{\rank}{\operatorname{rank}}
\newcommand{\cV}{\mathcal V}
\newcommand{\cR}{\mathcal R}
\newcommand{\cA}{\mathcal A}
\newcommand{\cL}{\mathcal L}
\newcommand{\cK}{\mathcal K}
\newcommand{\rev}{\operatorname{rev}}
\newcommand{\lean}[1]{\texttt{#1}}

\begin{document}

\title{An irrationality measure for Catalan's constant}

\author{Sean Litsey}
\email{sean.litsey.math@gmail.com}

\date{November 1, 2026}

\subjclass[2020]{Primary 11J82; Secondary 11J72, 15A15, 31A15, 68V20}

\keywords{Catalan's constant, irrationality measure, determinants, rank of reductions modulo p,
logarithmic potentials, Pick matrices, Heine's formula, Gauss quadrature, formal verification}

\begin{abstract}
Let $G=\sum_{j\ge0}(-1)^j(2j+1)^{-2}$ be Catalan's constant. We prove that $|G-a/b|>b^{-410}$ for
every integer $a$ and every sufficiently large integer $b$. Thus the irrationality exponent of $G$ is
at most $410$. The proof builds on OpenAI's proof that $G$ is irrational. That proof compares two
bounds for a determinant whose entries are linear forms in $1$ and $G$. The lower bound comes from
the denominators, and the upper bound from an integral representation. We improve both sides and add
the uniformity a measure needs. At the finite places we bound the rank of the reduction modulo $p$ of
the matrix of rational parts. This lowers the denominator at every prime in an explicit window. At the
real place we bound the determinant of the row functions by a signed estimate. A Pick matrix argument
reduces it to the block of the negative nodes. Heine's formula writes that block as a ratio of two
log-gas integrals, and both are bounded by logarithmic energy. A Gauss quadrature formula extends the
bound to an exponentially small neighbourhood of $G$, at a cost paid column by column. All steps are
formalized in the Lean proof assistant, with only the standard axioms.
\end{abstract}

\maketitle

\tableofcontents

\section{Introduction}\label{sec:intro}

\subsection{The result}
Catalan's constant
\[
 G=\sum_{j=0}^\infty\frac{(-1)^j}{(2j+1)^2}=L(2,\chi_{-4})=0.91596\,55941\,77219\ldots
\]
is the value at $s=2$ of the Dirichlet $L$-function of the nontrivial character modulo $4$. Its
arithmetic has long been studied by analogy with $\zeta(2)$ and $\zeta(3)$. Until recently it was not
known to be irrational. In October 2026 OpenAI \cite{OpenAI2026} proved that $G$ is irrational and
verified the proof in Lean.

The \emph{irrationality exponent} $\mu(\xi)$ of a real number $\xi$ is the infimum of the real
numbers $\mu$ such that $|\xi-a/b|>b^{-\mu}$ for all integers $a$ and all sufficiently large
integers $b$. Every irrational number has $\mu(\xi)\ge2$, and almost every real number has
$\mu(\xi)=2$. An \emph{irrationality measure} is an explicit upper bound for $\mu(\xi)$.

\begin{theorem}\label{thm:main}
For every integer $a$ and every sufficiently large integer $b$,
\[
 \left|G-\frac ab\right|>b^{-410}.
\]
In particular $\mu(G)\le410$.
\end{theorem}

The proof uses the prime number theorem. The theorem and every result it depends on have been verified
in the Lean~4 proof assistant on top of OpenAI's formalization. The formal proof has no unproved steps
and uses only the standard axioms (Section~\ref{sec:formal}). The use of AI in this work is disclosed at
the end of the paper.

\subsection{Background}
Irrationality proofs for values of $L$-functions at integers almost always follow the pattern of
Ap\'ery's proof for $\zeta(3)$ \cite{Apery79}. One constructs integer linear forms $q_n\xi-p_n$ that
tend to zero, with $q_n\ne0$, faster than the growth of their coefficients would allow for a rational
$\xi$. If $|q_n\xi-p_n|=e^{-\alpha n+o(n)}$ and $|q_n|=e^{\beta n+o(n)}$, then $\xi$ is irrational as
soon as $\alpha>0$, and $\mu(\xi)\le1+\beta/\alpha$. Beukers \cite{Beukers79} gave the linear forms for
$\zeta(2)$ and $\zeta(3)$ as integrals of Legendre-type polynomials. Hata \cite{Hata93} refined such
integrals. He used systematically the primes that divide all the coefficients of the forms, an idea that
goes back to Chudnovsky. Rhin and Viola \cite{RhinViola96,RhinViola01} found that families of
integrals for $\zeta(2)$ and $\zeta(3)$ are permuted by a finite group. The invariance of the integrals
lets one divide all their coefficients by a large common factor. This \emph{arithmetic method} gave
the sharpest measures for $\zeta(2)$ and $\zeta(3)$ for many years. Marcovecchio \cite{Marcovecchio09}
applied it to $\log2$. Zeilberger and Zudilin \cite{ZeilbergerZudilin20} obtained
$\mu(\pi)\le7.1032\ldots$ from Salikhov's integrals \cite{Salikhov08} with an additional arithmetic
saving.

For Catalan's constant the analogous linear forms exist but do not decay fast enough. Rivoal and Zudilin
\cite{RivoalZudilin2003} proved that one of the numbers $\beta(2),\beta(4),\ldots,\beta(2m)$ is
irrational for suitable $m$, where $\beta$ is the Dirichlet beta function. This does not single out
$G=\beta(2)$. Zudilin \cite{Zudilin2003} found an Ap\'ery-like recurrence for $G$. Its solutions give
linear forms that tend to zero, but too slowly to prove irrationality. Related constructions and their
limits are described in \cite{Zudilin02Cheb,KrattRivoal08,KrattZudilin19,Zudilin19Survey}. Eskandari,
Murty and Nemoto \cite{EMN2025} construct a mixed motive with period $G$ and obtain linear forms in $1$
and $G$ from it. Eskandari \cite{Eskandari2026} improves the approximation exponent of these forms. Leibtag
et al.\ \cite{LBKWSK26} found further Ap\'ery-like approximations to $G$ in conservative matrix fields.
Numerically the best of them still has a negative approximation exponent.

OpenAI's argument \cite{OpenAI2026} has a different shape. It uses a determinant $\Delta$ of size $48N$
instead of a sequence of linear forms. The entries of $\Delta$ are rational linear combinations of $1$
and $G$, given by double integrals against the Cauchy kernel $1/(1-ts)$. If $G$ were rational, the
denominators of the entries would give a lower bound $|\Delta|\ge e^{-(\mathsf A+o(1))n^2}$ with $n=48N$.
An integral representation gives the upper bound $|\Delta|\le e^{-(\mathsf U+o(1))n^2}$. It rests on
Andr\'eief's identity \cite{Forrester2018Andreief}, familiar from random matrix theory, and an energy
estimate shows $\mathsf U>\mathsf A$. The margin is small, about $1.3\cdot10^{-4}$. The nonvanishing of
$\Delta$ along the primes $N=p$ is proved separately by reduction modulo $p$. Determinants of linear forms
were used for irrationality before, in Zudilin's determinantal criterion with Hankel determinants
\cite{Zudilin}.

The first irrationality measure for $G$ was obtained by Jha \cite{Jha2026}, who proved
$\mu(G)\le10^8$. Carreira \cite{Carreira2026} treated the five integer parameters of the construction
as free. He stated conditions under which OpenAI's argument goes through for general parameters and
obtained $\mu(G)<2294$. He also formulated the best real-place constant of the method as an equilibrium
problem for two measures. In \cite{Carreira2026b} he records attempts to carry the method to $\beta(4)$.
Bojnal \cite{Bojnal2026} instantiated Carreira's general theorem at the parameters $(36,6,2,3,2)$ and
obtained $\mu(G)<2401.74$. The search and certification were carried out with Claude Fable~5.1. Xu
\cite{Xu2026} reported $\mu(G)\le727.03$ for the parameters $(20,3,1,2,1)$, with an improved trial
potential in the energy estimate. His certificate is numerical, in ball arithmetic, and is not
formalized. Davis \cite{Davis2026} obtained $\mu(G)\le528.23288$ with the design of this paper and the
arithmetic estimates of its formalization. The gain comes from new quadrature nodes. They are the
solutions of $T_K(P(t))=0$ for a fixed rational polynomial $P$ of degree $130$, on a support trimmed on
one side. This result is formalized in Lean. Arkhipov \cite{Arkhipov2026} modified that development and
reported $\mu(G)\le528.2327$. These node constructions are independent of the new ideas below, and we
have not combined them with ours. Sun's preprint \cite{Sun2026} had claimed a proof of the irrationality of $G$. Wachs
\cite{Wachs2026} reports that its main estimate is not justified.

Since September 2026 a series of new results in this area has appeared, several of them obtained with AI
systems. OpenAI proved $\mu(\pi)=2$, with a Lean formalization \cite{OpenAIPi2026}. Liu, Jiang and
Zhang proved $\mu(\log r)=2$ for every rational $r>0$ with $r\ne1$ \cite{LiuJiangZhang26}. Giraudin
lowered the bound for $\mu(\pi^2)$ to $5.0193784$ with two hypergeometric constructions of Zudilin
\cite{Zudilin14} at new parameters \cite{Giraudin26}. He also lowered the bound for $\mu(\zeta(3))$ to
$5.5138800$ \cite{Giraudin26z3}. Leibtag, Barkan, Kalisch, Weinbaum, Shalyt and Kaminer \cite{LBKWSK26}
showed that the classical constructions are one-dimensional trajectories in conservative matrix fields.
Searching over other trajectories, they obtained $\mu(\pi^2)<4.62675$ and $\mu(\zeta(3))<5.4197$.
Prellberg \cite{Prellberg26} proved that $\mathrm{Li}_2(-1/4)$, $\mathrm{Li}_2(1/5)$ and
$\mathrm{Li}_2(-1/3)$ are irrational. Fauzan announced a proof that $\zeta(5)$ is irrational, using
Hankel determinants \cite{Fauzan2026}. Lai combined this method with simultaneous approximations to show
that $1$, $\zeta(3)$ and $\zeta(4)$ are linearly independent over $\Q$ \cite{Lai2026}. For
$L(2,\chi_{-3})$, Calegari, Dimitrov and Tang proved that $1$, $\pi^2$ and $L(2,\chi_{-3})$ are linearly
independent over $\Q$ \cite{CDT2024}. Their effective holonomy bounds give $\mu(L(2,\chi_{-3}))\le24781$
\cite[Thm.~5.1]{CDT2025}.

\subsection{The companion paper}
Catalan's constant is a difference of two values of the trigamma function $\psi'$. For $0<s<t\le1$ put
\[
 \Psi(s,t)=\psi'(s)-\psi'(t)=\sum_{n=0}^\infty\Bigl(\frac1{(n+s)^2}-\frac1{(n+t)^2}\Bigr).
\]
Then $G=\frac1{16}\Psi(\frac14,\frac34)$, and also $L(2,\chi_{-3})=\frac19\Psi(\frac13,\frac23)$ and
$\pi^2=3\Psi(\frac12,1)$. The irrationality exponent does not change under $\xi\mapsto r\xi+r'$ with
$r,r'\in\Q$ and $r\ne0$. So a measure for $G$ is a measure for $\Psi(\frac14,\frac34)$. The companion
paper \cite{Litsey2026b} proves irrationality measures for the values $\Psi(s,1)$ with $s\in\{\frac14,
\frac34,\frac13,\frac23,\frac12\}$ and $s=\frac k5$, $\frac k7$, and for $\Psi(\frac13,\frac23)$. Among
them are $\mu(G\pm5\pi^2/48)$ and $\mu(L(2,\chi_{-3}))$. Its determinants are L\"owner matrices of digamma
differences. That method covers the pair $(\frac14,\frac34)$ as well, but its bound for this pair is too
weak to give a measure. The present paper therefore refines OpenAI's construction instead. The two papers
share their notation, their literature and the formalization repository.

\subsection{What a measure requires}\label{sec:intro-need}
Write $\Delta_N(X)$ for the determinant regarded as a polynomial in the constant, so that $\Delta_N(G)$
is OpenAI's determinant. Let $Q_N$ be the integer polynomial obtained from $\Delta_N$ by clearing
denominators. Three properties of $Q_N$ turn OpenAI's comparison into a measure.
\begin{enumerate}[label=(\roman*)]
\item \emph{Smallness on a neighbourhood of $G$.} The upper bound must hold not only at $X=G$. It must
hold for all $X$ with $|X-G|\le e^{-\rho N}$, and the exponent $\rho$ enters the final bound linearly.
\item \emph{Uniform nonvanishing.} One needs $Q_p(a/b)\ne0$ for primes $p$ with $p^2\nmid b$, in place of
$Q_p(G)\ne0$.
\item \emph{Arithmetic.} The integer $b^{\deg Q_N}Q_N(a/b)$ is nonzero, so it is at least $1$ in absolute
value. The smaller the denominator of $\Delta_N$ and its degree, the stronger the conclusion. The prime
$p$ must be chosen not to divide $b$, and a naive choice costs a factor.
\end{enumerate}
Let $d$ be the degree of $Q_N$ divided by $N$ and $\eta$ the decay rate of $Q_N$ normalised by $N^2$.
Then the conclusion is $\mu(G)\le d\rho/\eta$ (Proposition~\ref{prop:criterion}). The work is to make $d$
and $\rho$ small and $\eta$ large. Since $\eta$ is a small difference of two large numbers, the gains in
$\eta$ matter most.

\subsection{New ideas}\label{sec:intro-ideas}
\begin{enumerate}[label=\arabic*.,wide]
\item \emph{All parameters at once.} The determinant has five integer parameters
$(n_0,b_0,q_0,g_0,h_0)$, which we call a \emph{design}. Every argument below is carried out for an
arbitrary design, with the parameters as symbols. This covers the degree, the denominators at all
primes, the rank deficiencies, the nonvanishing and the estimates at the real place. A design is then
data. We use $(58,9,3,5,3)$, so that $n=58N$ and the degree is $53N$.

\item \emph{Rank of the reduction modulo $p$, in closed form.} Let $\nu_N(p)$ be the dimension of the left
kernel of the reduction modulo $p$ of $p^2$ times the matrix of rational parts. At a prime $p$ the
valuation of $\Delta_N$ gains $\nu_N(p)$. We prove four lower bounds for $\nu_N(p)$
(Theorem~\ref{thm:content}). Each is a piecewise linear function of $p$, with coefficients linear in the
five parameters, on its own window of primes. They come from antisymmetric combinations of rows, a
reflection of Chebyshev polynomials modulo $p$, and the Frobenius relation $(1-X)^p=1-X^p$. All four are
proved from one digit formula for the entries modulo $p$.

\item \emph{A counting inequality for the nullity.} The four proofs share one shape. A space $\cV$ of row
combinations is given, with a few linear conditions on $\cV$, and the conditions force all column
responses to vanish. We prove this once (Theorem~\ref{thm:count}). Then $\nu_N(p)\ge\dim\cV-\rank\Phi$,
where $\Phi$ is the map of conditions. A congruence from Lucas's theorem lets the inequality apply when
the column filter is longer than the prime. Further families of kernel vectors are instances of it.

\item \emph{A fixed rational subspace with stored kernel vectors.} In one form of the inequality the
space is a tensor product of polynomial coordinates with a fixed subspace $V\subset\Q^s$
(Theorem~\ref{thm:profile}). The required ranks are certified by integer vectors in the kernels of fixed
matrices, stored once. This treats $34$ windows of primes between $3N$ and $16.2N$.

\item \emph{Paying for the perturbation column by column.} The smallness near $G$ comes from a Gauss-type
quadrature formula for the coefficient of $G$, with $53N$ nodes. A perturbed determinant is a sum of
determinants in which some columns are replaced by point values. We bound each replaced column separately.
A binomial bound for generalized Vandermonde determinants (Lemma~\ref{lem:binomvand}) and a contour bound
for the weight of each column (Lemma~\ref{lem:keyhole}) give costs that peak at different columns. Their
sum peaks lower than the sum of their peaks.

\item \emph{A signed bound for the determinant of the row functions.} OpenAI's integrand contains a
determinant $\det\mathsf L$ that the energy method bounds only by $6^n$. We bound it on the configurations
that matter (Section~\ref{sec:signed}). A Pick matrix argument and a Schur complement reduce
$|\det\mathsf L|$ to the block of the negative nodes, at a cost of $11$ per positive node. The key input is
a contour bound for a quadratic form on the model space of a Blaschke product (Lemma~\ref{lem:T}).
Heine's formula and the Christoffel formula write the negative block as a ratio of two log-gas
integrals. We bound the numerator from above by completing squares in Chebyshev moments, and the
denominator from below by Jensen's inequality. An inequality for the logarithmic energy of signed
measures (Theorem~\ref{thm:reflection}) keeps the combined bound concave. A separate argument treats
positive nodes close to $1$.

\item \emph{Near-optimal configurations.} The signed bound needs hypotheses on the configuration. They hold
on every configuration whose energy is within the gain of the maximum. On the other configurations OpenAI's
bound is already strong enough. So the improved bound holds everywhere.

\item \emph{Trial potentials with complex poles, interpolated in the number of fixed columns.} The energy
estimates use trial sequences whose generating functions are logarithmic potentials of point charges near
the edges of the support of the extremal measure. The signed bound uses three trials, at fractions
$0$, $1/200$ and $1/16$ of fixed columns. Between them the fields are affine in the fraction, and the
remaining constant is an explicit quadratic.

\item \emph{A criterion with a free lower end of the window.} In the choice of the prime, the window of
candidates is $(V,U]$ with $U=\kappa\log b$ and $V=(\kappa-\omega)\log b$. We prove the criterion for every
$\omega\in(0,1]$ with the hypothesis $\kappa>2\omega$ (Proposition~\ref{prop:criterion}). We use
$\omega=3/5$.
\end{enumerate}

\subsection{The constants}
For the design $(58,9,3,5,3)$ the degree of $Q_N$ is $d=53N$. The perturbation exponent is
$\rho=\Rho$. The decay rate at the real place is $\TDelta\ldots$, and the denominators cost
$\mathsf A=\Aden\ldots$, after a gain of $\int\kappa\ge\ChiLean$ from the rank deficiencies. The signed
bound contributes $\Gain$ to the decay rate. The criterion uses $\eta=\EtaTot$, $\omega=3/5$,
$\delta=2.9\cdot10^{-7}$ and $\kappa=410/\rho=\Kappa\ldots$ Its last condition reads $B>\Bx\ldots$, and
$410$ is the least integer above it. Decimals are truncated. The constants of the formal proof are the
exact rational numbers.

\subsection{Organization and notation}
Section~\ref{sec:overview} sets up designs and states the rank deficiencies, the valuations, the
denominators, the bound at the real place and the criterion. Section~\ref{sec:setup} recalls OpenAI's
construction for a design and proves the degree, the denominators and the nonvanishing.
Section~\ref{sec:content} proves the rank deficiencies. Section~\ref{sec:real} reduces the bound near
$G$ to determinants with fixed columns and recalls the energy method. Section~\ref{sec:signed} proves the
signed bound. Section~\ref{sec:proof} proves Theorem~\ref{thm:main}. Section~\ref{sec:remarks} makes
remarks on the method, and Section~\ref{sec:formal} describes the formalization.

We use the notation of OpenAI's paper \cite{OpenAI2026} for the construction and refer to it for details.
Citations of the form \cite[Prop.~2.1]{OpenAI2026} point to its numbered statements. We write $\log$ for
the natural logarithm, $v_p$ for the $p$-adic valuation and $\theta(x)=\sum_{p\le x}\log p$ for
Chebyshev's function. $T_k$ and $U_k$ are the Chebyshev polynomials of the first and second kind, with
$U_{-1}=0$, $T_{-k}=T_k$ and $U_{-k-1}=-U_{k-1}$. We write $[X^k]F$ or $F[k]$ for the coefficient of $X^k$
in a polynomial $F$, $y_+=\max(y,0)$, and $\operatorname{rev}_EF(X)=X^EF(1/X)$ for the reversal of a
polynomial of degree at most $E$.

\section{Designs, rank deficiencies and the exponent}\label{sec:overview}

\subsection{Designs}
A \emph{design} is a quintuple $(n_0,b_0,q_0,g_0,h_0)$ of integers with $n_0,b_0,q_0,h_0\ge1$,
$g_0\ge0$ and $2g_0<n_0$. For $N\ge1$ we write
\begin{gather*}
 n=n_0N,\quad b=b_0N,\quad q=q_0N,\quad g=g_0N,\quad h=h_0N,\\
 L=n+b+q,\quad C=L+g,\quad H=C+h,\quad A=b+q+2g ,
\end{gather*}
and $d_0=n_0-g_0$. OpenAI's choice is the design $(48,7,4,4,2)$. We work throughout with
\begin{gather*}
 D=(58,9,3,5,3):\quad n=58N,\quad b=9N,\quad q=3N,\quad g=5N,\quad h=3N,\\
 L=70N,\quad C=75N,\quad H=78N,\quad A=22N,
\end{gather*}
and $d_0=53$. Section~\ref{sec:setup} recalls the construction of OpenAI's determinant for a design. It
yields an $n\times n$ matrix $\cA_N=\cA^0_N+G\cL_N$ with rational matrices $\cA^0_N$ and $\cL_N$. The
polynomial $\Delta_N(X)=\det(\cA^0_N+X\cL_N)$ has degree at most $n-g$ (Lemma~\ref{lem:degree}). The search
that led to $D$ is not part of the proof.

\subsection{Rank deficiencies}
For an odd prime $p$ with $H<p^2$ the entries of $p^2\cA^0_N$ are $p$-integral (Lemma~\ref{lem:digits}).
Let $\cA^0_{N,p}$ be its reduction modulo $p$, and let $\nu_N(p)$ be the dimension of its left kernel, that
is, $n$ minus its rank. Write $\Lambda=g+q+h$ and $e=(p+1)/2-q-h$.

\begin{theorem}\label{thm:content}
Let $D$ be a design, $N\ge1$, and $p$ an odd prime with $H<p^2$. Then $\nu_N(p)$ is at least each of the
following quantities whenever the stated conditions hold. Negative values are replaced by $0$.
\begin{enumerate}[label=(\alph*)]
\item (Theorem~\ref{thm:T}) If $2g<p$, $H\le4p$ and $D_T=\min(g,\,n-1-g-2p,\,C-3p)\ge1$:
\[
 D_T-(n+q+g+h-3p)_+ .
\]
\item (Theorem~\ref{thm:U}) If $2h\le p+1$, $L\le3p\le H$, $n+q+g+h\le3p$ and $2p\le n+q$:
\[
 \min(n,\,2p+2g+1)-g-2p-1 .
\]
\item (Theorem~\ref{thm:U6}) If $e\ge0$, $n\le4p+2g+1$, $L\le6p$ and $b+q\le p+\mu_0$, with
$\mu_0=\max(q,H+q-6p)$:
\[
 \min\Bigl(\frac{p-1}2-g,\ n+b+q+h-\frac{11p+1}2\Bigr).
\]
\item (Theorem~\ref{thm:S}) If $e\ge0$ and $n\le4p+2g+1$: $m_1-m_0-e$, with
\[
 m_0=\max(b+q,\,L+\Lambda-4p),\qquad m_1=\min(p,\,L-3p,\,b+q+2g+1).
\]
\end{enumerate}
\end{theorem}

These are the Lean theorems \lean{t\_rank}, \lean{u\_rank}, \lean{u6\_rank\_closed} and \lean{s\_rank},
proved for every admissible design. For $D$ they apply for $12N<p\le26N$, where they give most of the gain
(Section~\ref{ssec:kappa}).

Below $12N$ the closed formulas give nothing. Further kernel vectors come from a single counting
inequality, of which the proofs of Theorem~\ref{thm:content} are special cases. For $c\in\F_p^n$ and a raw
column $j\in[b,L)$ let $R_c(j)$ be the reduction modulo $p$ of $p^2\sum_rc_rR(r,j)$, with $R$ as in
\eqref{eq:rational-part}. These are the responses of Lemma~\ref{lem:digits}.

\begin{theorem}\label{thm:count}
Let $D$ be a design, $N\ge1$, $p$ an odd prime with $H<p^2$, and write $q=ap+r$ with $0\le r<p$. Let
$\cV$ be a finite-dimensional vector space over $\F_p$, $\iota\colon\cV\to\F_p^n$ an injective linear map and
$\Phi$ a linear map on $\cV$. Suppose that for every $x\in\ker\Phi$ there is a polynomial $G_x\in\F_p[X]$ of
degree less than $r$ with
\[
 \sum_{v=0}^a(-1)^{a-v}\binom av\,R_{\iota(x)}(j+vp)=G_x(j)\qquad(b\le j<L-ap).
\]
Then $\nu_N(p)\ge\dim\cV-\operatorname{rank}\Phi$.
\end{theorem}

For $q<p$ the hypothesis says that the conditions $\Phi$ force all raw responses to vanish. The case used
for the counting is the \emph{profile form}. There the space is a tensor product of a space of polynomial
coordinates with a fixed rational subspace, and the conditions are jets composed with fixed linear maps.

\begin{theorem}\label{thm:profile}
Let $v_1,\dots,v_d\in\Z^s$ span a subspace $V\subset\Q^s$ of dimension $d$. For $j$ in a finite set $J$ let
$B_j$ be an integer matrix with $s$ columns and $r_j\le d$ an integer such that there are $d-r_j$ linearly
independent vectors $k\in\Z^d$ with $B_j\bigl(\sum_ak_av_a\bigr)=0$. Let $\mathcal E\ne0$ be the product of a
nonzero $d\times d$ minor of $(v_1,\dots,v_d)$ and, for each $j$, of a nonzero maximal minor of these kernel
vectors. Let $p\nmid\mathcal E$ be an odd prime with $H<p^2$, write $q=ap+r$ as above, let $V_p=\F_p^d$ be the
space of coordinates of $V$ modulo $p$, and let $m$ and $w_j$ ($j\in J$) be integers $\ge0$. Suppose that there
are linear maps $\varphi_{j,k}\colon V_p^m\to V_p$ ($k<w_j$) and $\iota\colon V_p^m\to\F_p^n$ such that, on the joint
kernel of the maps $B_j\circ\varphi_{j,k}$, the map $\iota$ is injective and satisfies the hypothesis of
Theorem~\ref{thm:count}. Then
\[
 \nu_N(p)\ge m\,d-\sum_{j\in J}w_j\,r_j .
\]
\end{theorem}

In the applications $m$ and the $w_j$ are affine functions of $N$ and $p$ on a window of primes, while $V$
and the $B_j$ are fixed. So the bound is linear on the window. The data to be stored are the vectors $v_a$,
the matrices $B_j$ and the kernel vectors. For the design $D$ these bounds, the closed formulas and the
further families of Section~\ref{sec:content} combine into the following.

\begin{theorem}\label{thm:profile-D}
For the design $D$ there are $102$ intervals $(a_i,b_i]$, with $16/7=a_1<b_1\le a_2<\dots<b_{102}=53/2$, and a
function $\kappa$, linear on each interval and $0$ outside their union, with the following property. For all
large $N$ and every prime $p$ with $a_iN<p\le b_iN$ for some $i$,
\[
 \nu_N(p)\ge N\kappa(p/N)-2 .
\]
On each interval the bound is given by one of the families of Section~\ref{sec:content}. The area of
$\kappa$ is a rational number,
\[
 \int_0^\infty\kappa(x)\,dx=\frac{36235515874601938601596598}{1642609663271115992578125}=\Chi\ldots\ \ge\ \ChiLean .
\]
\end{theorem}

The contributions by ranges of $x=p/N$ are as follows. The third column counts the intervals on which the
bound is that of Theorem~\ref{thm:profile} with stored data.
\begin{center}
\begin{tabular}{@{}lrrl@{}}
\toprule
range of $x$ & intervals & with stored data & area \\
\midrule
$(16/7,\,3]$ & $11$ & $0$ & $0.06259$ \\
$(3,\,6]$ & $46$ & $23$ & $0.25136$ \\
$(6,\,12]$ & $27$ & $9$ & $0.71965$ \\
$(12,\,53/2]$ & $18$ & $2$ & $21.02610$ \\
\midrule
total & $102$ & $34$ & $\Chi$ \\
\bottomrule
\end{tabular}
\end{center}
The intervals in $(12,53/2]$ are listed in Section~\ref{ssec:kappa}. The full list is part of the
formalization. On each interval one bound is used. The families come from different arguments. Except for one
pair (Section~\ref{ssec:families}) their kernels are not shown to be independent, so at a prime the bounds are
not added.

\subsection{Valuations, denominators and the exponent}
The rank deficiencies enter through the following statement.

\begin{proposition}\label{prop:valuation}
Let $p$ be an odd prime with $H<p^2$ and let $z\in\Q$ have denominator prime to $p$. Then
$p^2(\cA^0_N+z\cL_N)$ has $p$-integral entries, its reduction modulo $p$ equals $\cA^0_{N,p}$, and
\[
 v_p\bigl(\Delta_N(z)\bigr)\ge\nu_N(p)-2n .
\]
\end{proposition}

At the other primes there are the valuation bounds of OpenAI's method (Section~\ref{sec:setup}). At $2$
there is a bound from the expansion into minors. Together they give the denominators.

\begin{proposition}\label{prop:denominators}
For the design $D$ there are positive integers $D_N$ with the following properties. $D_N\Delta_N(i)\in\Z$ for
every integer $i$, and $q^{2n}\mid D_N$ for every prime $N<q\le2N$. For every $\varepsilon>0$ and all large $N$
\begin{gather*}
 \log D_N\le(\mathsf A+\varepsilon)N^2,\\
 \mathsf A=\frac{2673}8\log2+116+\frac{24269}4-\int_0^\infty\kappa(x)\,dx=\Aden\ldots
\end{gather*}
\end{proposition}

Here $\frac{2673}8\log2$ is the contribution of the prime $2$. The number $116+\frac{24269}4$ is the area of
the piecewise linear function that bounds the poles at the odd primes (Section~\ref{sec:setup}). The last
term is the gain of Theorem~\ref{thm:profile-D}.

At the real place we prove the following in Sections~\ref{sec:real} and~\ref{sec:signed}. Put
\begin{equation}\label{eq:rho-def}
 \rho=n_0c_2+\Psi_0=\Rho,\qquad c_2=\Ctwo,\quad \Psi_0=\PsiZero .
\end{equation}
Here $c_2$ is the energy cost of one fixed column and $\Psi_0$ bounds the logarithmic size of the weight
and the binomial factor of one column, per unit of $N$.

\begin{theorem}\label{thm:real}
Let $U_0=U+\Gain/n_0^2=\Uzero\ldots$, with $U=\Utarget$ and $n_0=58$. For every $\varepsilon>0$, all large
$N$ and all real $x$ with $|x-G|\le e^{-\rho N}$,
\[
 |\Delta_N(x)|\le\exp\bigl(-(n_0^2(U_0-\tfrac12\log2)-\varepsilon)N^2\bigr).
\]
\end{theorem}

The number $n_0^2(U-\frac12\log2)=\TDeltaOld\ldots$ is the decay rate that OpenAI's energy method gives
for the design $D$. The signed bound adds $\Gain$, so $n_0^2(U_0-\frac12\log2)=\TDelta\ldots$ Since
$\mathsf A=\Aden\ldots$, the polynomial $Q_N=(d_0N)!\,D_N\Delta_N$ has integer coefficients, degree at
most $d_0N$, and
\[
 |Q_N(x)|\le e^{-\eta N^2}\qquad(|x-G|\le e^{-\rho N}),\qquad \eta=\EtaTot .
\]
The exact difference is $\EtaExact\ldots$, and $\eta$ is a rational number just below it. The remaining
properties of Section~\ref{sec:intro-need} are the nonvanishing of $Q_p(a/b)$ for $p^2\nmid b$
(Proposition~\ref{prop:nonvanishing}) and the divisibility of the $i$-th coefficient of $Q_N$ by $q^{2i}$
for $N<q\le2N$ (Lemma~\ref{lem:recycling}). They are used to choose the prime. The criterion is the
following.

\begin{proposition}\label{prop:criterion}
Let $\xi\in\R$, $d\in\N$, and let $\eta,\rho,\kappa,\delta,\omega>0$ satisfy
\[
 \omega\le1,\quad 2\omega<\kappa,\quad \eta<2d,\quad 1<2\bigl(\omega-\delta(2\kappa-\omega)\bigr),\quad
 d+4d\delta\kappa<\eta\kappa .
\]
Let $Q_p\in\Z[X]$ ($p$ prime) satisfy the following for all large $p$. $\deg Q_p\le dp$. $Q_p(a/b)\ne0$
whenever $a\in\Z$, $b\ge1$ and $p^2\nmid b$. For all primes $p<q\le2p$ and $i\le dp$, $q^{2i}$ divides the
$i$-th coefficient of $Q_p$. Finally $|Q_p(x)|\le e^{-\eta p^2}$ for $|x-\xi|\le e^{-\rho p}$. Then
\[
 |\xi-a/b|>b^{-\rho\kappa}
\]
for all integers $a$ and all sufficiently large $b$.
\end{proposition}

The conditions say, in the order of the proof in Section~\ref{sec:proof}, three things. The window of primes
contains a prime not dividing $b$. The primes of the window above the chosen one are recycled. The final
inequality is contradictory. For the design $D$ one takes $d=53$, $\omega=3/5$, $\delta=2.9\cdot10^{-7}$
and $\kappa=B/\rho$ with $B=410$. The last condition then reads $B>x:=d\rho/(\eta-4d\delta)=\Bx\ldots$

\subsection{A reading of the exponent}
For a fixed design the bound is $d\rho/\eta$ in the limit. Its three ingredients come from three
different places. The degree $d=n-g$ comes from the construction. The perturbation exponent $\rho$ comes
from the quadrature and the energy cost of a column. The decay rate $\eta$ is the difference of the real
place and the denominators, $\TDelta$ and $\Aden$. Both large numbers are improved here: the rank
deficiencies lower the second by $22.05$, and the signed bound raises the first by $\Gain$. Without either,
$\eta$ would be $\EtaInf$.

\section{OpenAI's construction for a design}\label{sec:setup}

We recall the construction of \cite{OpenAI2026} in the generality in which we use it, and we prove for an
arbitrary design the degree bound, the bounds for the denominators at the finite places and the
nonvanishing at the primes $N=p$. Carreira \cite{Carreira2026} checks that the argument of OpenAI goes
through for designs satisfying a few conditions. We need the same for the finite places and, in addition,
for the entries modulo $p^2$.

\subsection{The determinant}
Let $D=(n_0,b_0,q_0,g_0,h_0)$ be a design, with the notation of Section~\ref{sec:overview}. For $0\le r<n$
the row polynomials are
\begin{equation}\label{eq:rows}
 P_r(t)=(1-t)^ht^{C-1}T_{|r-g|}(1/t),\qquad \deg P_r<H,
\end{equation}
together with companion polynomials $D_r$ built from $U_{|r-g|-1}$ \cite[(3)]{OpenAI2026}. Thus
$P_r$ depends only on $|r-g|$. Put $f(t)=\sqrt{1-t^2}$ and $d\mu(t)=|t|\,dt/f(t)$ on $(-1,1)$, and
\[
 A_r(t)=P_r(t)-\frac32\,\mathbf 1_{\{t>0\}}\frac{f(t)}{t}D_r(t),\qquad
 \psi_k(s)=s^{b+k}(1-s)^q .
\]
The determinant studied by OpenAI is that of the $n\times n$ matrix
\begin{equation}\label{eq:matrix}
 \cA_N(r,k)=\int_{-1}^1\int_0^1 A_r(t)\,\frac{\psi_k(s)}{1-ts}\,ds\,d\mu(t)
 \qquad(0\le r,k<n).
\end{equation}
The factor $(1-s)^q$ in $\psi_k$ amounts to a finite difference of order $q$ applied to the columns.
The entries are rational linear combinations of $1$ and $G$: there are rational matrices $\cA^0_N$
and $\cL_N$ with
\begin{equation}\label{eq:L}
 \cA_N=\cA^0_N+G\,\cL_N,\qquad
 \cL_N(r,k)=\sum_{v=0}^{q}(-1)^v\binom qv\sum_{i}[t^i]P_r\cdot4\,c_{(i-b-k-v)/2},
\end{equation}
where $c_l=4^{-l}\binom{2l}{l}$ for integers $l\ge0$ and $c_z=0$ otherwise
\cite[Prop.~2.1 and (20)]{OpenAI2026}. The rational part is
\begin{equation}\label{eq:rational-part}
 \begin{gathered}
 \cA^0_N(r,k)=\sum_{v=0}^q(-1)^v\binom qv R(r,b+k+v),\\
 R(r,j)=\sum_{i<H}P_r[i]\,M(i,j)-\frac32\sum_{i<H}D_r[i]\,Z(i,j)
 \end{gathered}
\end{equation}
with rational numbers $M(i,j)$ (moments of $\mu$ against the Cauchy kernel) and $Z(i,j)$, where
$Z(i,i)=-\sum_{k\le i}k^{-2}$ and $Z(i,j)=(H_i-H_j)/(i-j)$ for $i\ne j$, $H_i=\sum_{k\le i}1/k$
\cite[Prop.~2.1]{OpenAI2026}. The design enters \eqref{eq:rational-part} only through the integer
coefficients $P_r[i],D_r[i]$ and the range of $j$, which is $[b,L)$. Put
$\Delta_N(X)=\det(\cA^0_N+X\cL_N)\in\Q[X]$, so that $\Delta_N(G)=\det\cA_N$.

\subsection{The degree}
\begin{lemma}\label{lem:degree}
$\deg\Delta_N\le n-g$.
\end{lemma}

\begin{proof}
By \eqref{eq:L} the row $r$ of $\cL_N$ depends only on $P_r$, and hence only on $|r-g|$. So the rows
$r$ and $2g-r$ of $\cL_N$ agree for $0\le r<g$, and $2g-r\le n-1$ because $2g<n$. Subtracting row $r$ from
row $2g-r$ for each $r<g$ does not change the determinant and makes $g$ rows of the new coefficient of
$X$ vanish. A determinant $\det(B+XL')$ in which $L'$ has $n-g$ nonzero rows has degree at most $n-g$ in $X$.
\end{proof}

\subsection{Rank deficiency and valuation}
We use repeatedly the following elementary fact.

\begin{lemma}\label{lem:rank-valuation}
Let $M$ be an $n\times n$ matrix with entries in $\Z_{(p)}$ whose reduction modulo $p$ has rank at most
$n-e$. Then $v_p(\det M)\ge e$.
\end{lemma}

\begin{proof}
Factor the reduction through $\F_p^{n-e}$ and lift the two factors to integral matrices $U,V$, so that
$M=UV+pE$ with $E$ integral at $p$. Expanding the determinant multilinearly in the columns, every term
that uses more than $n-e$ columns of $UV$ vanishes, and every other term has at least $e$ factors $p$.
\end{proof}

\subsection{The digit formula}
The reduction of $p^2\cA^0_N$ modulo $p$ is described by the digits in base $p$ of the coefficients of the
row polynomials. This is the one place where the arithmetic of the moments enters. It is design-free.
For an odd prime $p$ and $c\in\F_p^n$ put
\[
 P_c=\sum_rc_rP_r,\qquad D_c=\sum_rc_rD_r,\qquad W=(1-X^2)^{(p-1)/2},
\]
(reading the $P_r,D_r$ as polynomials in $X$), and let $F[k]$ denote the coefficient of $X^k$.

\begin{lemma}[digit formula]\label{lem:digits}
Let $p$ be an odd prime with $H<p^2$, let $0\le j<H$, and write $j=Jp+\ell$ with $0\le\ell<p$. Then
$p^2R(r,j)$ is $p$-integral for every row $r$, and for every $c\in\Z^n$
\[
 p^2\sum_rc_rR(r,j)\equiv\sum_{u\ge0}\Bigl(\alpha_{u,J}\,(XWP_c)[up+\ell]+\beta_{u,J}\,D_c[up+\ell]\Bigr)
 \pmod p ,
\]
where $\alpha_{u,J}=M(u-1,J)$ and $\beta_{u,J}=-\tfrac32Z(u,J)$ are $p$-integral numbers that do not depend on
$c$, on $\ell$ or on the design.
\end{lemma}

\begin{proof}
Two facts about the numbers $M$ and $Z$ are needed, for $0\le i,j<p^2$. Write $j=Jp+\ell$ and let
$\varrho\in[0,p)$ be the residue of $j-i-1$ modulo $p$, $u=(i+1+\varrho-\ell)/p$. First, $p^2M(i,j)$ is
$p$-integral and $p^2M(i,j)\equiv W[\varrho]\,M(u-1,J)\pmod p$. This is the digit reduction of the moments in
\cite[\S3]{OpenAI2026}, obtained from the $p$-adic analysis of their recursions, and it does not involve a design.
Second, $p^2Z(i,j)$ is $p$-integral and is congruent to $Z(\lfloor i/p\rfloor,J)$ if $i\equiv j\pmod p$ and to
$0$ otherwise. For the second, write $i'=\lfloor i/p\rfloor$. Since $i<p^2$, $H_i=p^{-1}H_{i'}+S_i$ and
$H^{(2)}_i=p^{-2}H^{(2)}_{i'}+S'_i$ with $S_i,S'_i$ $p$-integral, which gives the case $i=j$. If $i\not\equiv j$,
then $p^2Z(i,j)=p\cdot pZ(i,j)\equiv0$. If $i\equiv j$, $i\ne j$, then $i-j=(i'-J)p$, and the integers
$k\in(j,i]$ prime to $p$ form $i'-J$ complete systems of residues modulo $p$, so that
$\sum_k1/k\equiv(i'-J)\sum_{r=1}^{p-1}1/r\equiv0\pmod p$. Hence $H_i-H_j\equiv(H_{i'}-H_J)/p$ and
$p^2Z(i,j)\equiv Z(i',J)$.

Now use \eqref{eq:rational-part}. By the first fact the part $\sum_iF[i]\,p^2M(i,j)$ with $F=P_c$ is
$\sum_iF[i]\,W[\varrho(i)]\,M(u(i)-1,J)$. For fixed $u$ the indices with $u(i)=u$ are those for which
$\varrho=up+\ell-1-i$ lies in $[0,p-1]$, and $W$ has degree $p-1$, so the sum over these $i$ is
$(XFW)[up+\ell]$. This gives the $\alpha$-terms. By the second fact the part with $D_c$ is
$\sum_{i\equiv\ell}D_c[i]\,Z(\lfloor i/p\rfloor,J)=\sum_uZ(u,J)D_c[up+\ell]$, which gives the $\beta$-terms.
The $p$-integrality of $M(u-1,J)$ and $Z(u,J)$ for $u,J<p$ is part of the first and second facts.
\end{proof}

For the filtered columns we use the following form. Let $0\le k<n$, $s=b+k$, $J=\lfloor s/p\rfloor$ and
$\mu=(s\bmod p)+q$. If $\mu<p$, the $q+1$ raw columns $s,\dots,s+q$ that enter column $k$ have the same
digit $J$ and residues $(s\bmod p)+v<p$. Since
$\sum_{v=0}^q(-1)^v\binom qvF[t+v]=[X^{t+q}]\bigl((X-1)^qF\bigr)$ for every polynomial $F$,
Lemma~\ref{lem:digits} gives
\begin{equation}\label{eq:filtered-digits}
 p^2\sum_rc_r\cA^0_N(r,k)\equiv\sum_{u\ge0}\Bigl(\alpha_{u,J}\,Q_c[up+\mu]+\beta_{u,J}\,\widehat D_c[up+\mu]\Bigr),
\end{equation}
where $Q_c=(X-1)^q\,XWP_c$ and $\widehat D_c=(X-1)^qD_c$. A combination $c$ whose responses vanish on all
raw columns in $[b,L)$ lies in the left kernel of $\cA^0_{N,p}$, and a combination whose responses on the
columns outside a set vanish and which kills the set by \eqref{eq:filtered-digits} does as well.

\subsection{Integrality and the valuation at $p$}
\begin{proof}[Proof of Proposition~\ref{prop:valuation}]
By Lemma~\ref{lem:digits} the entries of $p^2\cA^0_N$ are $p$-integral, since the filter is an integer
combination of raw columns with $j<L\le H$. The entries of $\cL_N$ are $p$-integral for odd $p$, since
$P_r[i]\in\Z$ and $c_l\in\Z[1/2]$, so $p^2\cL_N\equiv0\pmod{p^2}$ and, for $z$ with denominator prime to
$p$, $p^2(\cA^0_N+z\cL_N)$ has $p$-integral entries with the same reduction as $p^2\cA^0_N$. By
Lemma~\ref{lem:rank-valuation} the determinant of the latter has valuation at least $\nu_N(p)$, and
$\det\bigl(p^2(\cA^0_N+z\cL_N)\bigr)=p^{2n}\Delta_N(z)$.
\end{proof}

\subsection{The denominators}
Let $p$ be an odd prime with $H<p^2$ and $p\le H$. The method of \cite[\S3]{OpenAI2026}, applied to a
design, bounds from below the valuation of $\Delta_N(z)$ for every $z$ with denominator prime to $p$ by
$-\ell_N(p)$, where $\ell_N(p)$ counts the columns whose entries have a pole at $p$, once the coefficients of
the row polynomials are grouped by their digits in base $p$ and the column filter is taken into account. It
is given by two formulas, in which all differences are truncated at $0$. For $2p\le H$,
\[
 \ell_N(p)=2n-\Bigl(\min\bigl(p,\,L-p,\,A\bigr)-\max\bigl(b,\,H-2p\bigr)-q\Bigr),
\]
and for $H<2p\le2H$, with $R=(L-b)-(\min(A,L-p)-b)-(\min(p,L)-\max(b,H-p))$,
$P'=\min(A,L-p)-b$ and $E'=\min(b,H-p)-(L-p)$,
\[
 \ell_N(p)=\min\bigl(2n,\ n+R,\ 2R+P'+E'\bigr).
\]
These are the exponents of the generic valuation theorems in the formalization. For the primes that
interest us they are the same as in OpenAI's paper for its design. For the design $D$ the function
$\ell_N(p)=Nf(p/N)$ is piecewise linear in $x=p/N$ ($\lfloor\cdot\rfloor$ rounding aside):
\[
 f(x)=\begin{cases}
 116&x\le59/2,\\
 175-2x&59/2\le x\le69/2,\\
 106&69/2\le x\le39,\\
 184-2x&39\le x\le81/2,\\
 265-4x&81/2\le x\le48,\\
 217-3x&48\le x\le69,\\
 148-2x&69\le x\le70,\\
 78-x&70\le x\le78,
 \end{cases}
\]
and $f=0$ for $x\ge78$, since for $p>H$ the matrix $\cA^0_N+z\cL_N$ is $p$-integral. Its integral over
$[1,78]$ is $24269/4$, and $f=116=2\cdot58$ on $[0,1]$. Note that $f(x)=2n_0$ for $x\le59/2$, which
contains all the primes in $(N,2N]$ and, in particular, the whole range $(16/7,53/2]$ of the rank
deficiencies of Theorem~\ref{thm:profile-D}. There the bound of Proposition~\ref{prop:valuation} improves
$-\ell_N(p)=-2n$ to $\nu_N(p)-2n$.

The primes $p\le\sqrt H$ are not covered by these bounds. For them we use the cruder bound
$6n\log_pH$ for the exponent, which contributes $O(\sqrt N\,N\log N)=o(N^2)$ in total.

\subsection{The prime $2$}
At $p=2$ the entries are not integral and the bound comes from the expansion of the determinant into
minors, in the manner of \cite[\S3]{OpenAI2026}. Its exponent is governed by a quadratic form in the
numbers $m,\ell$ of columns of two exceptional types, which is
\[
 \tfrac32m^2+\tfrac12\ell^2-(H+C-b)\,m-C\,\ell+Cn ,
\]
and the following identity shows that it is at least $-\mathsf A_2$ whenever $m,\ell\ge0$,
$m+\ell\le n$ and $C\ge n$, where
\begin{equation}\label{eq:A2}
 \mathsf A_2=\frac{(n+H-b)^2}8-\frac{n^2}2 .
\end{equation}
With $s=n-m-l$ and $m_*=(n+H-b)/4$,
\[
 \tfrac32m^2+\tfrac12\ell^2-(H+C-b)m-C\ell+Cn+\mathsf A_2
 =2(m-m_*)^2+s\,(C-n+m)+\tfrac12s^2\ge0 .
\]
This gives a $2$-adic bound $v_2(\Delta_N(i))\ge-\mathsf A_2-O(N\log N)$, uniform in the integer $i$. For $D$,
$\mathsf A_2=(58+78-9)^2/8-58^2/2=2673/8$ in units of $N^2$.

\subsection{The common denominator}
\begin{proof}[Proof of Proposition~\ref{prop:denominators}]
Let $\kappa_N(p)$ be the largest of the lower bounds for $\nu_N(p)$ of Section~\ref{sec:content} that apply at
$p$ (and $0$ if none does), for primes $p\le H$ with $H<p^2$, and put $e_N(p)=\min\bigl(\ell_N(p),\,2n-\kappa_N(p)\bigr)$. For the
primes with $p^2\le H$ use the crude bound above, for $p=2$ the $2$-adic bound, and let $D_N$ be the
product of $p^{e_N(p)}$ over $p\le H$. By Proposition~\ref{prop:valuation} and the bounds of this section,
$v_p(\Delta_N(i))\ge-e_N(p)$ for every integer $i\ge0$ and every $p$ (for $p>H$ the matrix is
$p$-integral), so $D_N\Delta_N(i)\in\Z$. A prime $q\in(N,2N]$ has $\kappa_N(q)=0$, since no window
contains it, and $\ell_N(q)=2n$, so $q^{2n}\mid D_N$ (here $N\ge79$ guarantees $q>\sqrt H$).

For the size of $D_N$, the sum of $\ell_N(p)\log p$ over the odd primes is, by the prime number theorem in
the weighted form
\begin{equation}\label{eq:weighted-pnt}
 \frac1N\sum_{p}F(p/N)\log p\longrightarrow\int_0^\infty F(x)\,dx\qquad(N\to\infty),
\end{equation}
for bounded, Riemann integrable, compactly supported $F$, equal to $N^2\int_0^{78}f(x)\,dx+o(N^2)$ with
$\int_0^{78}f=116+24269/4$. The primes in the intervals of Theorem~\ref{thm:profile-D} lie in
$(16N/7,53N/2]$ and have $\ell_N(p)=116N$, so their exponent is at most $116N-N\kappa(p/N)+2$. The sum of
$(N\kappa(p/N)-2)\log p$ over these primes is $N^2\int\kappa+o(N^2)$ by \eqref{eq:weighted-pnt}. The $2$-adic part contributes
$\mathsf A_2\log2\,N^2+o(N^2)$.
\end{proof}

The polynomial $\Delta_N$ has degree at most $n-g=d_0N$ and takes values in $D_N^{-1}\Z$ at the
integers. An integer-valued polynomial of degree $m$ has coefficients in $\frac1{m!}\Z$ (Newton's
interpolation formula), so
\begin{equation}\label{eq:QN}
 Q_N=(d_0N)!\,D_N\,\Delta_N(X)\in\Z[X],\qquad \deg Q_N\le d_0N,
\end{equation}
and $\log(d_0N)!=o(N^2)$. In this way the gain from the rank deficiencies enters through the integer
polynomial itself, and not through a separate common factor of its coefficients.

\begin{lemma}\label{lem:recycling}
For all large $N$ and every prime $N<q\le2N$, the $j$-th coefficient of $Q_N$ is divisible by $q^{2j}$.
\end{lemma}

\begin{proof}
For these primes the matrices $q^2\cA^0_N$ and $\cL_N$ are $q$-integral by Lemma~\ref{lem:digits}, and
$q^{2n}$ divides $D_N$. The identity
\[
 Q_N(X/q^2)=(d_0N)!\,\frac{D_N}{q^{2n}}\det\bigl(q^2\cA^0_N+X\cL_N\bigr)
\]
shows that $Q_N(X/q^2)$ has $q$-integral coefficients. Its $j$-th coefficient is $c_{N,j}/q^{2j}$.
\end{proof}

\subsection{Nonvanishing}
\begin{proposition}\label{prop:nonvanishing}
Let $D$ be a design and suppose that the three $n_0\times n_0$ matrices $B_\sigma$, $\sigma\in\{0,1,-1\}$,
defined below, are invertible. Then for all large primes $p$ and every rational $z$ such that $pz$ has
denominator prime to $p$, $\Delta_p(z)\ne0$. In particular $\Delta_p(a/b)\ne0$ whenever $a\in\Z$, $b\ge1$
and $p^2\nmid b$.
\end{proposition}

The matrices are those of OpenAI's reduction modulo $p$ at $N=p$, with the design in place of $48$: let
$B$ be the rational $(n_0+1)\times n_0$ matrix of the entries of the filtered rational part at $N=1$ for
the rows $0,\dots,n_0$ (the row $n_0$ is an auxiliary row), and let $B_\sigma$ have rows
$B_r+\sigma B_{r+1}$, $0\le r<n_0$. For the design $D$, $B$ has $59$ rows and $58$ columns.

\begin{proof}
Let $p>4H_0$, with $H_0=H/N=78$, so that $H=H_0p<p^2/4$, and let $p$ be prime to the numerators and denominators
of the entries of $B$ and of the three determinants. In the construction at $N=p$ write the row index
$r=pr_0+i$ with $0\le r_0<n_0$ and $0\le i<p$. Since $q=q_0p$, $C=C_0p$, $h=h_0p$, $g=g_0p$, the Frobenius
relation $(1-t)^{p}\equiv1-t^p$ expresses the row polynomials modulo $p$ through those of the design at
$N=1$ and the successor row $r_0+1$. This is where the auxiliary row enters. It is valid for rows on
both sides of $g$ because the polynomials are defined through $T_{r-g}$ and $U_{r-g-1}$ for signed
indices. The filter satisfies $s^{b_0p+k}(1-s)^{q_0p}=s^{\ell}(s^p)^{b_0+k_0}(1-s^p)^{q_0}$ in $\F_p[s]$, so
the filtered columns of one residue class use the raw columns $b_0+k_0,\dots,b_0+k_0+q_0$ of the design at
$N=1$ and never change the residue. By \cite[Lemma~3.1 and eqs.~(31), (35), (36)]{OpenAI2026}, which
need only $H<p^2/4$, the reduction of $p^2\cA^0_p$ modulo $p$ is then, after ordering the rows by $(i,r_0)$ and
the columns by $(\ell,k_0)$, the matrix $a^{\mathsf T}\otimes B_0+b^{\mathsf T}\otimes B_1$, where $a,b$ are the
fixed $p\times p$ matrices of OpenAI's identity (51), with $\det a=2^{-p(p-1)/2}$. The permutation
matrix $\Pi$ in this decomposition has the eigenvalue $0$ once and $\pm1$ each $(p-1)/2$ times, and
therefore
\[
 \det=(\det a)^{n_0}\,\det B_0\;\det(B_0+B_1)^{(p-1)/2}\det(B_0-B_1)^{(p-1)/2}\ne0
\]
in $\F_p$. Hence $p^{2n}\Delta_p(0)$ is a $p$-adic unit. For $z$ with $pz$ of denominator prime to $p$, the
difference $p^2z\cL_p=p\,(pz)\,\cL_p$ is $p$-integral with zero reduction (the entries of $\cL_p$ are
$p$-integral), so the reduction of $p^2(\cA^0_p+z\cL_p)$ is the same invertible matrix and
$v_p(\Delta_p(z))=-2n$. If $p^2\nmid b$, then $p(a/b)$ has denominator prime to $p$.
\end{proof}

For the design $D$ the invertibility of $B_0$, $B_0+B_1$ and $B_0-B_1$ is certified by exhibiting their
inverses modulo the prime $1009$, which exceeds $H_0=78$ so that every denominator of $B$ is a unit
there. The three determinants are then nonzero rational numbers. The computation modulo $1009$ certifies
the base matrices. The conclusion for all large primes is a theorem and not an extrapolation from
finitely many primes. Nothing in the proof uses the proportions of the design except $2g<n$, and $H<p^2/4$.

\section{Corank modulo $p$}\label{sec:content}

This section proves the gain in the denominators. The principle is that of the arithmetic method of Hata
and Rhin--Viola, adapted to a determinant. If the reduction modulo $p$ of $p^2\cA^0_N$ has corank $e$, then
$\Delta_N(z)$ gains a factor $p^e$ at every $z$ with denominator prime to $p$
(Proposition~\ref{prop:valuation}). So we need explicit row combinations $c\in\F_p^n$ with
$c^{\mathsf T}\cA^0_{N,p}=0$.

Here is the heuristic. Modulo $p$ the Chebyshev polynomials of index $kp+d$ factor through $X^p$, so rows
placed symmetrically about the centres $g+kp$ are related. Suitable combinations of such rows have row
polynomials with long runs of vanishing base-$p$ digits. By the digit formula their values at most moment
columns then vanish, and the few remaining values are removed by imposing linear conditions. Each surviving
dimension is a unit of corank.

We first prove the counting inequality (Theorem~\ref{thm:count}). It says that a space of row combinations
subject to a few linear conditions lies in the left kernel. We then derive four closed formulas from it,
which give most of the gain. Finally we describe the further families and the resulting function $\kappa$.
Three facts about characteristic $p$ are used throughout: the Chebyshev identities
\begin{equation}\label{eq:cheb-p}
 \begin{gathered}
 T_p=X^p,\qquad U_{p-1}=(X^2-1)^{(p-1)/2},\\
 T_{kp}=T_k(X^p),\qquad U_{kp-1}=(X^2-1)^{(p-1)/2}\,U_{k-1}(X^p)
 \end{gathered}
\end{equation}
in $\F_p[X]$, the relation $(1-X)^p=1-X^p$, and the digit formula \eqref{eq:column-digits}. The last two
identities of \eqref{eq:cheb-p} follow from $T_{kp}=T_k\circ T_p$ and $U_{kp-1}=(U_{k-1}\circ T_p)\,U_{p-1}$.

\subsection{Setting}
Fix parameters, an integer $N\ge1$ and an odd prime $p$ with $H<p^2$. For $0\le r<n$ put $m_r=r-g$, so that
$-g\le m_r\le n-1-g$ and $n-1-g\ge g$. In the notation of \eqref{eq:rows},
\[
 P_r=(1-X)^h\operatorname{rev}_{C-1}T_{m_r},\qquad D_r=(1-X)^h\operatorname{rev}_{C-1}U_{m_r-1},
\]
where $T_{-m}=T_m$ and $U_{-m-1}=-U_{m-1}$, so that $D_g=0$. For $c\in\F_p^n$ write
$S_c=\sum_rc_r\operatorname{rev}_{C-1}T_{m_r}$ and $U_c=\sum_rc_r\operatorname{rev}_{C-1}U_{m_r-1}$, so that
$P_c=(1-X)^hS_c$ and $D_c=(1-X)^hU_c$. The lowest degree of $\operatorname{rev}_{C-1}T_m$ is at least
$C-1-|m|\ge A=b+q+2g$, and that of $\operatorname{rev}_{C-1}U_{m-1}$ is at least $A+1$. Put $\chi=g+q+h$. Then
the polynomials $Q_c$ and $\widehat D_c$ of \eqref{eq:Qc} have their coefficients in $[A+1,\,L+\chi+p-1]$
and $[A+1,\,L+\chi-1]$ respectively.

A vector $c$ lies in the left kernel of $\cA^0_{N,p}$ as soon as the right side of
\eqref{eq:column-digits} vanishes for every column $k$. We only use columns whose $q+1$ moment columns
share the same base-$p$ digit.

\subsection{The counting inequality}
All the kernels of this section arise in the same way. A space of row combinations is chosen, a few linear
conditions are imposed, and the conditions make every value vanish. The dimension of the space minus the
number of conditions is then a lower bound for the corank. The following theorem states this once. It also
allows $q\ge p$.

\begin{theorem}\label{thm:count}
Let $N\ge1$ and let $p$ be an odd prime with $H<p^2$. Write $q=ap+r$ with $0\le r<p$. Let $\cV$ be a
finite-dimensional vector space over $\F_p$, $\iota\colon\cV\to\F_p^n$ an injective linear map and $\Phi$ a
linear map on $\cV$. Suppose that for every $x\in\ker\Phi$ there is a polynomial $g_x\in\F_p[X]$ of degree
less than $r$ with
\[
 \sum_{v=0}^a(-1)^{a-v}\binom av\,R_{\iota(x)}(j+vp)=g_x(j)\qquad(b\le j<L-ap).
\]
Then $\nu_N(p)\ge\dim\cV-\operatorname{rank}\Phi$.
\end{theorem}

For $q<p$ the hypothesis says that the conditions $\Phi$ force all values at the moment columns to vanish.

\begin{proof}
By \eqref{eq:rational-part}, the entry of $c^{\mathsf T}\cA^0_{N,p}$ at column $k$ is
$\sum_{v=0}^q(-1)^v\binom qvR_c(b+k+v)$. This is the polynomial $(1-T)^q$ applied to $R_c$, where $T$ acts as
the shift. Modulo $p$ we have $(1-T)^p=1-T^p$, hence
\[
 (1-T)^q=(1-T)^r(1-T^p)^a .
\]
This is Lucas's congruence for the binomial coefficients $\binom qv$. So the entry at $k$ equals
$\sum_{u=0}^r(-1)^u\binom ru\,\Sigma_c(b+k+u)$, where
$\Sigma_c(j)=\sum_{v=0}^a(-1)^v\binom av R_c(j+vp)$. The arguments satisfy $b\le j$ and
$j+ap\le b+n-1+q=L-1$. For $c=\iota(x)$ with $x\in\ker\Phi$ we have $\Sigma_c(j)=\pm g_x(j)$. The difference
of order $r$ of a polynomial of degree less than $r$ vanishes identically. So $\iota(\ker\Phi)$ lies in the
left kernel of $\cA^0_{N,p}$. Since $\iota$ is injective, $\nu_N(p)\ge\dim\ker\Phi=\dim\cV-\operatorname{rank}\Phi$.
\end{proof}

Suppose instead that the conditions only force the values on a set $\mathcal C$ of columns to vanish. The
same argument, with the remaining $n-|\mathcal C|$ entries added to the conditions, gives
$\nu_N(p)\ge\dim\cV-\operatorname{rank}\Phi-(n-|\mathcal C|)$. The proofs of the four closed formulas below
are of this form. In Theorems~\ref{thm:T} and~\ref{thm:U} the set $\mathcal C$ consists of all columns, and in
Theorems~\ref{thm:U6} and~\ref{thm:S} it is a block of columns.

\subsection{Four closed formulas}
Each closed formula exhibits explicit vectors in the left kernel. The rows used are symmetric about the
centres $g+kp$, and the identities \eqref{eq:cheb-p} make their combinations vanish modulo $p$. Recall
$e=(p+1)/2-q-h$.

\begin{theorem}\label{thm:content}
Let $N\ge1$ and let $p$ be an odd prime with $H<p^2$. Then $\nu_N(p)$ is at least each of the following
quantities whenever the stated conditions hold. Negative values are replaced by $0$.
\begin{enumerate}[label=(\alph*)]
\item (Theorem~\ref{thm:T}) If $2g<p$, $H\le4p$ and $k_T=\min(g,\,n-1-g-2p,\,C-3p)\ge1$:
\[
 k_T-(n+q+g+h-3p)_+ .
\]
\item (Theorem~\ref{thm:U}) If $2h\le p+1$, $L\le3p\le H$, $n+q+g+h\le3p$ and $2p\le n+q$:
\[
 \min(n,\,2p+2g+1)-g-2p-1 .
\]
\item (Theorem~\ref{thm:U6}) If $e\ge0$, $n\le4p+2g+1$, $L\le6p$ and $b+q\le p+\varsigma_0$, with
$\varsigma_0=\max(q,H+q-6p)$:
\[
 \min\Bigl(\frac{p-1}2-g,\ n+b+q+h-\frac{11p+1}2\Bigr).
\]
\item (Theorem~\ref{thm:S}) If $e\ge0$ and $n\le4p+2g+1$, then for every integer $k\ge2$ the quantity $m_1-m_0-e$, with
\[
 m_0=\max(b+q,\,L+\chi-kp),\qquad m_1=\min(p,\,L-(k-1)p,\,A+1).
\]
\end{enumerate}
\end{theorem}

For our parameters these bounds apply for $\frac{74}5N<p\le\frac{954}{29}N$, where they give most of the gain
(Section~\ref{ssec:kappa}). We prove the four parts in turn.

\subsection{Antisymmetric combinations of rows}
\begin{theorem}\label{thm:T}
Let $2g<p$, $H\le4p$ and $k_T=\min(g,\,n-1-g-2p,\,C-3p)\ge1$. Then
$\nu_N(p)\ge k_T-(n+q+g+h-3p)_+$.
\end{theorem}

The rows used are the differences of rows placed symmetrically about the three centres $g$, $g+p$ and
$g+2p$.

\begin{proof}
For $1\le d\le k_T$ and $k=0,1,2$ let $\epsilon_k(d)=e_{g+kp+d}-e_{g+kp-d}\in\F_p^n$. These rows exist
because $d\le g$ and $g+2p+d\le n-1$. The $3k_T$ vectors $\epsilon_k(d)$ are linearly independent, since their
supports are disjoint. Indeed the supports lie in the intervals $[kp,kp+2g]$ of row indices, which are
disjoint because $2g<p$. The row $g+kp\pm d$ has $m=kp\pm d$. By $T_{a+b}-T_{a-b}=2(X^2-1)U_{a-1}U_{b-1}$
and $U_{a+b-1}-U_{a-b-1}=2T_aU_{b-1}$ together with \eqref{eq:cheb-p},
\begin{align*}
 T_{kp+d}-T_{kp-d}&=2(X^2-1)^{(p+1)/2}\,U_{k-1}(X^p)\,U_{d-1},\\
 U_{kp+d-1}-U_{kp-d-1}&=2\,T_k(X^p)\,U_{d-1}.
\end{align*}
Put $R_d=2(1-X)^hX^{C-d}\operatorname{rev}_{d-1}U_{d-1}$, the value of $D_{\epsilon_0(d)}$. Reverse at degree
$C-1$ and use $\operatorname{rev}_{kp}(T_k(X^p))=(\operatorname{rev}_kT_k)(X^p)$, with $T_0=1$, $T_1=X$ and
$T_2=2X^2-1$. One finds
\[
 D_{\epsilon_1(d)}=X^{-p}R_d,\qquad D_{\epsilon_2(d)}=2X^{-2p}R_d-R_d,
\]
\begin{gather*}
 XWP_{\epsilon_0(d)}=0,\qquad XWP_{\epsilon_1(d)}=X^{-p}R_d-X^pR_d,\\
 XWP_{\epsilon_2(d)}=2X^{-2p}R_d-2R_d .
\end{gather*}
The polynomial $R_d$ is supported in degrees $[C-d,H-1]$, which lies in $[3p,4p)$ because $d\le C-3p$ and
$H\le4p$. So at each residue $\ell$ it has a single base-$p$ digit $\varrho_\ell(d)$, in position $u=3$. For
$x=(x_0,x_1,x_2)\in\F_p^3$ the value of $\sum_kx_k\epsilon_k(d)$ at the moment column $Jp+\ell$ is therefore,
by Lemma~\ref{lem:digits}, $\varrho_\ell(d)\,f_J(x)$ with
\[
 f_J(x)=x_0\beta_{3,J}+x_1(\alpha_{2,J}-\alpha_{4,J}+\beta_{2,J})+x_2(2\alpha_{1,J}-2\alpha_{3,J}+2\beta_{1,J}-\beta_{3,J}).
\]
This depends neither on $d$ nor on $\ell$. The moment columns with $J\ge3$ do not meet the support, since
$3p+\ell\ge C-d\ge C-g=L$. Choose $x^*\ne0$ with $f_1(x^*)=f_2(x^*)=0$, which is two linear equations in three
unknowns. For $c=\sum_ds_d\sum_kx^*_k\epsilon_k(d)$ the values at the columns with $J=1,2$ vanish. The columns
with $J=0$ have $\ell\ge b$, and the value there is $f_0(x^*)\sum_ds_d\varrho_\ell(d)$, the coefficient of
$X^{3p+\ell}$ in $\sum_ds_dR_d$. These are at most $(H-3p-b)_+=(n+q+g+h-3p)_+$ linear conditions on
$s=(s_d)$. The map $s\mapsto c$ is injective, so the solutions form a space of dimension at least
$k_T-(n+q+g+h-3p)_+$ in the left kernel.
\end{proof}

\subsection{A reflection of Chebyshev polynomials}
The next bound rests on a reflection formula modulo $p$.

\begin{lemma}\label{lem:u3-reflect}
Let $p$ be an odd prime and $\varpi=(p+1)/2$. In $\F_p[X]$, for every integer $m$,
\[
 T_{p-m}=X^pT_m-(X^2-1)^{\varpi}\,U_{m-1}.
\]
\end{lemma}

\begin{proof}
Both sides satisfy the three-term recurrence in $m$, so it suffices to check $m=0$ and $m=1$. For $m=0$ it
is $T_p=X^p$, by \eqref{eq:cheb-p}. For $m=1$ use $T_{p+1}=XT_p+(X^2-1)U_{p-1}$ and $T_{p-1}=2XT_p-T_{p+1}$.
\end{proof}

\begin{theorem}\label{thm:U}
Let $2h\le p+1$, $L\le3p\le H$, $n+q+g+h\le3p$ and $2p\le n+q$, and put $n^*=\min(n,2p+2g+1)$. Then
$\nu_N(p)\ge n^*-g-2p-1$.
\end{theorem}

\begin{proof}
We restrict $c$ to the rows $r<n^*$, so that $m_r\le n^*-1-g\le2p+g$. Apply Lemma~\ref{lem:u3-reflect} with
$m=m_r$ and reflect at degree $p+C-1$. For every such row this gives
\[
 (1-X^2)^\varpi\operatorname{rev}_{C-1}U_{m-1}-X\operatorname{rev}_{C-1}T_m=-X^{p+1}\operatorname{rev}_{C-1}T_{p-m}.
\]
The right side is divisible by $X^{p+C-|p-m|}$, hence by $X^L$ since $|p-m|\le p+g=p+C-L$. Multiplying by
$(1-X)^h$ and summing, we get for every $c$ supported on these rows
\begin{equation}\label{eq:identityK}
 X^L\mid(1-X^2)^\varpi D_c-XP_c .
\end{equation}

Let $\cV$ be the space of such $c$ with $(1-X^2)^\varpi\mid P_c$. Since $(1-X)^h$ divides every $P_r$, this is
$\varpi-h$ conditions at $X=1$ and $\varpi$ conditions at $X=-1$, that is, at most $p+1-h$ linear
conditions. For $c\in\cV$ write $P_c=(1-X^2)^\varpi Y$ and $\widetilde Q=XY$. Since $(1-X^2)^p=1-X^{2p}$,
\[
 XWP_c=(1-X^{2p})\,\widetilde Q .
\]
The polynomial $P_c$ has no coefficients below $X^{A^*}$, $A^*=C-n^*+g$, and none from $X^H$ on. So
$\widetilde Q$ is supported in $[A^*+1,\,H-p-1]$. Dividing \eqref{eq:identityK} by the unit $(1-X^2)^\varpi$
modulo $X^L$ gives $D_c[k]=\widetilde Q[k]$ for $k<L$.

Now let $H-3p\le\ell<L-2p$. The inequality $L-2p\le A^*+1$ is equivalent to $n^*\le2p+2g+1$, and
$L-2p\le p$ since $L\le3p$. By the supports, $\widetilde Q[\ell]=0$, $\widetilde Q[2p+\ell]=0$ and
$\widetilde Q[3p+\ell]=0$, and $D_c[3p+\ell]=0$ because $3p+\ell\ge H$. So with
$\varkappa=\widetilde Q[p+\ell]$ the base-$p$ digits of $XWP_c$ at the residue $\ell$ are
$(0,\varkappa,0,-\varkappa,0,\dots)$, and those of $D_c$ are $(0,\varkappa,0,0,\dots)$. Hence, by
Lemma~\ref{lem:digits}, the value at every moment column $Jp+\ell$ with such an $\ell$ is $\varkappa$ times a
constant depending on $J$ only. On $\cV$ these columns, which have $J\in\{0,1,2\}$, therefore span a space
of dimension at most one for each residue $\ell$.

Put $\ell_0=H-3p$ and $\ell_1=L-2p$. The hypotheses give $\ell_0\le b\le\ell_1$. For $\ell_0\le\ell<b$ the
moment columns $j\in[b,L)$ with this residue have $J\in\{1,2\}$, and for $b\le\ell<\ell_1$ they have
$J\in\{0,1,2\}$. So the saving is $2\ell_1-\ell_0-b$. The number of linear conditions needed to make all
values on the $n+q=L-b$ moment columns vanish on $\cV$ is therefore
$(p+1-h)+(n+q)-(2\ell_1-\ell_0-b)$. Substituting $\ell_1=L-2p$ and $\ell_0=H-3p$, this is $2p+1+g$. So the
space of solutions in the span of $n^*$ rows has dimension at least $n^*-g-2p-1$.
\end{proof}

\subsection{A block of columns: the digits $1$ to $5$}
\begin{theorem}\label{thm:U6}
Let $e=(p+1)/2-q-h\ge0$, $n\le4p+2g+1$ and $L\le6p$, and put $\varsigma_0=\max(q,H+q-6p)$ and
$\varsigma_1=L-5p$. Assume $b+q\le p+\varsigma_0$. Then $\nu_N(p)\ge\varsigma_1-\varsigma_0-e$, that is,
\[
 \nu_N(p)\ge\min\Bigl(\frac{p-1}2-g,\ n+b+q+h-\frac{11p+1}2\Bigr).
\]
The first term is the minimum if $6p\le H$ and the second if $6p\ge H$.
\end{theorem}

In Theorems~\ref{thm:T} and~\ref{thm:U} the conditions are imposed on the rows. Here the argument concerns a
block of columns. We first need an identity.

\begin{lemma}\label{lem:K6}
In $\F_p[X]$, for every integer $k$, $2X^p(X^2-1)^{(p-1)/2}T_k-(2X^{2p}-1)U_{k-1}=U_{2p-k-1}$. If
$n\le4p+2g+1$, then for every $c$, $X^L\mid 2XWS_c-(2-X^{2p})U_c$, and consequently
\begin{equation}\label{eq:K6}
 2\,Q_c[i]=2\,\widehat D_c[i]-\widehat D_c[i-2p]\qquad(i<L).
\end{equation}
\end{lemma}

\begin{proof}
For all integers $m,k$ one has $U_{m-1}T_k-T_mU_{k-1}=U_{m-k-1}$, since both sides satisfy the same
recurrence in $k$ and agree for $k=0,1$. Take $m=2p$ and use \eqref{eq:cheb-p}: $T_{2p}=2X^{2p}-1$ and
$U_{2p-1}=2T_pU_{p-1}=2X^p(X^2-1)^{(p-1)/2}$. Reflect at degree $2p+C-1$. The first term becomes
$2XW\operatorname{rev}_{C-1}T_k$, the second $(2-X^{2p})\operatorname{rev}_{C-1}U_{k-1}$, and the right side
$X^{2p+C-1}U_{2p-k-1}(1/X)$. For $k=m_r\ge-g$ with $2p-k-1\ge0$ the right side is divisible by $X^{C+k}$, hence
by $X^L$. If $2p-k-1<0$, it has degree $k-2p-1$ and is divisible by $X^{4p+C-k}$, hence by $X^L$ when
$k\le4p+g$. This holds because $k\le n-1-g\le4p+g$. Summing with weights $c_r$ and multiplying by
$(X-1)^q(1-X)^h$ gives \eqref{eq:K6}.
\end{proof}

\begin{proof}[Proof of Theorem~\ref{thm:U6}]
Let $\mathcal M=[\varsigma_0,\varsigma_1)$ and $w=\varsigma_1-\varsigma_0$. We may assume $w>e$. The block consists of the columns
$k(J,\varsigma)=Jp+\varsigma-(b+q)$ for $1\le J\le5$ and $\varsigma\in \mathcal M$. These indices are nonnegative by the hypothesis
$b+q\le p+\varsigma_0$, and smaller than $n$ because $5p+\varsigma_1-1=L-1$. There are $5w$ of them. As $\varsigma<p$, formula
\eqref{eq:column-digits} applies, and the value of $c$ at $k(J,\varsigma)$ is
$\sum_u\alpha_{u,J}Q_c[up+\varsigma]+\beta_{u,J}\widehat D_c[up+\varsigma]$.

\emph{Supports.} For $\varsigma\in \mathcal M$ we have $p+\varsigma\le p+\varsigma_1-1=L-4p-1\le A$ by $n\le4p+2g+1$. So the digits $u=0,1$
of $Q_c$ and $\widehat D_c$ vanish. From $\varsigma_0\ge H+q-6p=L+\chi-6p$ we get that the digits $u\ge6$ of
$\widehat D_c$ and $u\ge7$ of $Q_c$ vanish. For $u\le5$ we have $up+\varsigma\le L-1$, so \eqref{eq:K6} applies.

\emph{From $Q$ to $\widehat D$.} For $u=2,\dots,5$, \eqref{eq:K6} gives
$\widehat D_c[up+\varsigma]=Q_c[up+\varsigma]+\tfrac12\widehat D_c[(u-2)p+\varsigma]$, and for $u=2,3$ the last index is at most $A$,
so that term vanishes. Hence if $Q_c[up+\varsigma]=0$ for $u=2,\dots,5$ and all $\varsigma\in \mathcal M$, then by induction
$\widehat D_c[up+\varsigma]=0$ for $u=2,\dots,5$, and all digits of $\widehat D_c$ in $\mathcal M$ vanish.

\emph{A relation from the Frobenius map.} We have
$Q_c=(-1)^qX(1-X)^{q+h+(p-1)/2}(1+X)^{(p-1)/2}S_c$, and $q+h+(p-1)/2+e=p$. Since $(1-X)^p=1-X^p$,
\[
 (1-X)^e\,Q_c=(1-X^p)\,Y_c,\qquad Y_c=(-1)^qX(1+X)^{(p-1)/2}S_c,
\]
and $Y_c$ is supported in $[A+1,\,C+(p-1)/2]$. Compare the coefficients at $up+\varsigma$, subtract the
coefficient at $(u-1)p+\varsigma$, and sum over $u=2,\dots,6$. The right side telescopes to
$Y_c[6p+\varsigma]-Y_c[p+\varsigma]$, which vanishes for $\varsigma\ge\varsigma_0+e$. The second term vanishes because $p+\varsigma\le A$,
and the first because $6p+\varsigma\ge H+q+e=C+(p+1)/2>\deg Y_c$. So, if $Q_c[up+\varsigma']=0$ for $u=2,\dots,5$ and all
$\varsigma'\in \mathcal M$, then
\[
 \sum_{t=0}^e(-1)^t\binom et\,Q_c[6p+\varsigma-t]=0\qquad(\varsigma_0+e\le\varsigma<\varsigma_1).
\]
This is a recurrence with unit leading coefficient. So if $Q_c[6p+\varsigma]=0$ for the first $e$ values
$\varsigma=\varsigma_0,\dots,\varsigma_0+e-1$, then it vanishes for all $\varsigma\in \mathcal M$.

\emph{Counting.} Impose on $c$ that the values at the $n-5w$ columns outside the block vanish, that
$Q_c[up+\varsigma]=0$ for $2\le u\le5$ and $\varsigma\in \mathcal M$ ($4w$ conditions), and that $Q_c[6p+\varsigma]=0$ for the first $e$
values of $\varsigma$. By the three steps above all the values at the block columns vanish. These are at most
$n-5w+4w+e=n-(w-e)$ linear conditions, so the left kernel has dimension at least $w-e$. Finally, $w-e$
equals $\varsigma_1-q-e=n+b+q+h-(11p+1)/2$ if $\varsigma_0=q$, and $(p-1)/2-g$ if $\varsigma_0=H+q-6p$.
\end{proof}

\subsection{The digits $1$ to $k$}\label{ssec:S}
The proof of Theorem~\ref{thm:U6} uses five digits. A shorter block uses fewer. The count is the same for
every number of digits $k\ge2$, and we prove it once. The cases $k=4$ and $k=3$ both occur for our
parameters. The case $k=3$ gives the bounds on $(\frac{186}7N,28N]$.

\begin{theorem}\label{thm:S}
Let $e=(p+1)/2-q-h\ge0$ and $n\le4p+2g+1$. Let $k\ge2$ be an integer and put
\[
 m_0=\max(b+q,\,L+\chi-kp),\qquad m_1=\min\bigl(p,\,L-(k-1)p,\,A+1\bigr).
\]
Then $\nu_N(p)\ge m_1-m_0-e$.
\end{theorem}

The proof rests on a statement about digits only. Write $\mathcal M=[m_0,m_1)$. For $\varsigma\in\mathcal M$ and
$u\ge0$ put $Q_u=Q_c[up+\varsigma]$ and $D_u=\widehat D_c[up+\varsigma]$, with the polynomials of
\eqref{eq:Qc}. The relation \eqref{eq:K6} reads $2Q_u=2D_u-D_{u-2}$ for $up+\varsigma<L$, with $D_{u-2}=0$
for $u<2$. The Frobenius relation $(1-X)^eQ_c=(1-X^p)Y_c$ of the proof of Theorem~\ref{thm:U6} gives
\begin{equation}\label{eq:frob-digits}
 \sum_{t=0}^e(-1)^t\binom et\,Q_c[i-t]=Y_c[i]-Y_c[i-p]\qquad(i\ge0).
\end{equation}

\begin{lemma}\label{lem:digit-vanishing}
Let $w=m_1-m_0>e$ and assume the hypotheses of Theorem~\ref{thm:S}. Suppose that $D_j=0$ for
$1\le j<k$ and all $\varsigma\in\mathcal M$, and that $Q_k=0$ for the first $e$ residues
$\varsigma=m_0,\dots,m_0+e-1$. Then $Q_u=D_u=0$ for all $u\ge0$ and all $\varsigma\in\mathcal M$.
\end{lemma}

\begin{proof}
Since $m_1\le A+1$, we have $\varsigma\le A$ and the digits $u=0$ of $Q_c$ and $\widehat D_c$ vanish. Since
$\varsigma\ge L+\chi-kp$, the digits $u\ge k$ of $\widehat D_c$ and $u\ge k+1$ of $Q_c$ vanish. Indeed these
polynomials have no coefficients beyond $L+\chi-1$ and $L+\chi+p-1$. The digits $D_j$ with $1\le j<k$ vanish by
assumption. For $1\le j<k$ we have $jp+\varsigma\le(k-1)p+m_1-1<L$. So \eqref{eq:K6} applies. It gives
$2Q_j=2D_j-D_{j-2}=0$, because $D_{j-2}$ is either absent or a digit already known to vanish. Only $Q_k$
remains.

Let $\varsigma\ge m_0+e$. Apply \eqref{eq:frob-digits} at $i=up+\varsigma$ for $u=0,\dots,k+1$ and add. The right sides
telescope to $Y_c[(k+1)p+\varsigma]$. This coefficient vanishes. Indeed $Y_c$ has degree at most
$C+(p-1)/2$, while $(k+1)p+\varsigma\ge p+L+\chi>C+(p-1)/2$, because $p+q+h>(p-1)/2$. On the left, every residue
$\varsigma-t$ with $0\le t\le e$ lies in $\mathcal M$. So all digits vanish except the digit $k$, and
\[
 \sum_{t=0}^e(-1)^t\binom et\,Q_c[kp+\varsigma-t]=0 .
\]
The term $t=0$ has coefficient $1$. The first $e$ values of $Q_k$ vanish by assumption. Induction on $\varsigma$
shows $Q_k=0$ on all of $\mathcal M$.
\end{proof}

\begin{proof}[Proof of Theorem~\ref{thm:S}]
We may assume $w=m_1-m_0>e$. The columns $Jp+\varsigma-(b+q)$ with $0\le J<k$ and $\varsigma\in\mathcal M$ exist.
They are nonnegative since $\varsigma\ge b+q$, and smaller than $n$ since $(k-1)p+\varsigma<L$. They are
distinct since $\varsigma<p$. So there are $kw$ of them. By \eqref{eq:column-digits} the value of $c$ at such a
column is a combination of the digits $Q_u,D_u$ at the residue $\varsigma$. Impose on $c$ that the values at the
other $n-kw$ columns vanish, that $D_j=0$ for $1\le j<k$ and all $\varsigma\in\mathcal M$ ($(k-1)w$ conditions),
and that $Q_k=0$ for the first $e$ residues. By Lemma~\ref{lem:digit-vanishing} all digits at the residues of
$\mathcal M$ vanish. So the values at the $kw$ columns vanish as well. There are at most
$(n-kw)+(k-1)w+e=n-(w-e)$ conditions, and the left kernel has dimension at least $w-e$.
\end{proof}

The lemma is the digit-only core of the argument. The proof uses no determinant and no lower bound for $N$.


\subsection{Decompositions modulo $p$}
The further families organize the row combinations as follows.

\begin{definition}\label{def:slot}
A \emph{decomposition modulo $p$} of a family of row combinations consists of
\begin{enumerate}[label=(\alph*)]
\item a linear map $x\mapsto c(x)$ from $(\F_p^K)^t$ to $\F_p^n$, injective on a subspace of admissible
 $x=(x_1,\dots,x_t)$;
\item polynomials $\pi_1,\dots,\pi_t\in\F_p[X]$ of degree less than $e_{\max}$;
\item linear maps $\mathcal T,\mathcal G\colon\F_p^K\to\F_p[Y]$,
\end{enumerate}
such that for admissible $x$
\begin{gather*}
 (X-1)^hP_{c(x)}=\sum_d\pi_d(X)\,\mathcal T(x_d)(X^p),\\
 (X-1)^h(X^2-1)^{\frac{p-1}2}D_{c(x)}=\sum_d\pi_d(X)\,\mathcal G(x_d)(X^p).
\end{gather*}
\end{definition}

The identities \eqref{eq:cheb-p} and $(1-X)^p=1-X^p$ provide such decompositions. By the digit formula, the
value of $c(x)$ at the moment column $Jp+\ell$ is a linear form in the $x_d$:
\[
 \sum_{0\le e'<e_{\max},\ p\mid H-1-\ell-e'}\ \Theta_{J,\,(H-1-\ell-e')/p}\Bigl(\sum_d\pi_d[e']\,x_d\Bigr).
\]
Here $\pi_d[e']$ is the coefficient of $X^{e'}$, and each $\Theta_{J,s}$ is a linear form in $J$ and the digit
$s$, determined by $\mathcal T$ and $\mathcal G$.

\subsection{Further families}\label{ssec:families}
The following lower bounds for $\nu_N(p)$ are proved for all parameters, by explicit choices of the space,
the conditions and the columns in Theorem~\ref{thm:count} or its variant. Each is a piecewise linear function
of $p$ and the parameters, positive on intervals described by linear inequalities. We describe them in
words. The statements are those of the formalization.
\begin{enumerate}[label=(F\arabic*),wide]
\item \emph{Frobenius products.} The row combinations have polynomials that are the product of a fixed
factor, a free polynomial of bounded degree $d_P$ in the Chebyshev variable, and a polynomial $\Xi(X^p)$ with
exponents in an interval of $K$ consecutive values. The conditions require $\Xi$ to annihilate the
functionals seen by the columns. If $\#\text{seen}$ of these functionals are independent, the bound is
$(d_P+1)(K-\#\text{seen})$. For our parameters this gives $66N-2p-1$ on $(\frac{97}3N,\frac{954}{29}N]$.
\item \emph{A truncated form of Theorem~\ref{thm:U6}.} Without the hypothesis $n\le4p+2g+1$, the bound
$\varsigma_1-\varsigma_0-e_+-(n-4p-2g-1)_+$ holds whenever $L\le6p$ and $b+q\le p+\varsigma_0$. For our parameters it gives
$5p-74N$ on $(\frac{74}5N,15N]$ and $p-14N$ on $(15N,16N]$, up to a bounded term.
\item \emph{A tail of the antisymmetric family.} With the four centres $g+kp$ ($0\le k\le3$) and
$k_T'=n-1-g-3p$, the bound is $p-1-(2g+q+h)$, under $4p\le L$, $b\le p$, $k_T'\le g$, $k_T'\le(p-1)/2$,
$3p+1+g\le n$ and $4p+b\le H$. For our parameters this is $p-20N-1$ on $(\frac{222}{11}N,\frac{146}7N]$.
\item \emph{Two families of decompositions and their sum.} The first fixes a digit $s\le H/p$ and a range of
blocks of rows, with polynomials $\pi_d$ divisible by a power of $X-1$. The second has polynomials
$(X^p-1)\pi$ with $\pi$ supported in an interval of degrees. Under a separation condition their kernels are
independent when both apply, at the same digit $s$ or at the digits $s$ and $s-1$. Then the sum of the two
bounds is a lower bound. This is the one place where two bounds are added. For our parameters it gives
bounds on six intervals between $8.4N$ and $10.9N$.
\item \emph{Residue intervals and rational cells.} For $p\le q$ one chooses the row combinations residue by
residue. On an interval $[\alpha,\beta)$ of residues of constant combinatorial type the bound is
$\beta-\alpha+p-q$ (Theorem~L$^*$). For $p\ge h$ and for $p\le h$ one uses cells defined over $\Q$. A cell
fixes a window of the digits, a range of blocks of rows and an integer profile, such that the value
of the combination at the columns of the cell has a fixed difference over $\Q$. For a prime $p$ prime to the modulus of the cell
the bound is $(\mathrm{hi}-\mathrm{lo})-(p-q)$ in the first case (Theorem~C) and
$(\mathrm{hi}-\mathrm{lo})-(h-p)-(p-q)$ in the second (Theorem~C$'$), where $[\mathrm{lo},\mathrm{hi})$ is an interval of positions
compatible with the cell. For our parameters each of these theorems supplies selected subintervals, with gaps,
of an interval. Theorem~C$'$ gives bounds on $(\frac{20}7N,\frac{66}{23}N]$ and $(\frac{80}{27}N,3N]$.
Theorem~L$^*$ gives bounds on $(\frac{74}{21}N,\frac{18}5N]$, $(\frac{37}{10}N,\frac{72}{19}N]$ and
$(\frac{74}{19}N,4N]$. Theorem~C gives bounds on $(4N,\frac{37}9N]$, $(\frac{72}{17}N,\frac{74}{17}N]$,
$(\frac92N,\frac{87}{19}N]$ and $(\frac{97}{21}N,\frac{37}8N]$. On $(\frac{87}{19}N,\frac{97}{21}N]$ the bound
combines two cells of Theorem~C and has the value $\mathrm{hi}_A+(\mathrm{hi}_B-\mathrm{lo}_B)-(p-4N)$, with
$\mathrm{lo}_B=97N-21p$, $\mathrm{hi}_A=\min(97N-21p,10N-2p)$ and $\mathrm{hi}_B=\min(10N-2p,70N-1-15p)$.
Intervals shifted by tiny rational amounts are listed in the formalization.
\item \emph{Affine cells, Theorem~HY$^+$ and Theorem~MT.} On $(\frac{83}{31}N,\frac{101}{37}N]$ and
$(\frac{36}{13}N,\frac{14}5N]$ the same
counting applies with cells on which the values of the combinations at consecutive digits differ by a fixed
slope, for primes $p>44$ prime to a fixed integer. Between these two intervals the bounds come from explicit
subspaces. Theorem~HY$^+$ uses the residue sums for the digits $s=11,10,9$ on the following intervals. For $s=11$
it gives $(\frac{70}9N,8N]$. For $s=10$ it gives $(\frac{17}2N,\frac{97}{11}N]$, and a variant with a short fold gives
$(\frac{97}{11}N,\frac{4409091}{500000}N]$. For $s=9$ it gives $(\frac{66}7N,\frac{202}{21}N]$, and a variant with a short
fold gives $(\frac{97}{10}N,\frac{8300000001}{850000000}N]$. A mixed family of antisymmetric and symmetric pairs
of rows (Theorem~MT) gives $\nu_N(p)\ge\frac{29}2p-77N-2$ on $(\frac{43}8N,\frac{97}{18}N]$.
\item \emph{Theorem Z.} Two further families, called the laws $\mathrm{L}$ and $\mathrm{H}$ in the formalization,
give bounds on $(\frac{186}{11}N,\frac{160}9N]$ and on $(\frac{132}7N,\frac{222}{11}N]$. Write $T=H-5p$ and $h_p=(p-1)/2$.
The law $\mathrm L$ gives
\begin{multline*}
 2\nu_N(p)\ge11p-2(n+b+q+g)-1\\
 {}+2\min(0,L-5p)+2\min(0,T-b)+2\min(0,h+h_p-T).
\end{multline*}
In the regime $5p\le L$ and $b\le T\le h+h_p$ the three minima vanish. For our parameters the three selected lines
are $\frac{11}2p-93N$, $\frac p2-7N$ and $80N-\frac92p$, up to a loss of at most two. The two laws count the conditions
that the Taylor coefficients of the combinations impose on a space of rows.
\end{enumerate}
The interval $(\frac{96666667}{10000000}N,\frac{290999999}{30000000}N]$ belongs to an explicit subspace of
Appendix~\ref{app:windows}, with $\kappa(\theta)=\frac{599999999}{100000000}-\frac\theta2$. It is not assigned to
Theorem~HY$^+$. The interval $(\frac{202}{21}N,\frac{29}3N]$ before it belongs to (F4).

In the cases that need a prime to avoid finitely many bad values, only $p$ larger than fixed integers occur,
and the corresponding $N$ are large.

For the primes below $12N$ that none of these closed forms covers, we use the counting inequality in a special
form. The space of row combinations is the tensor product of a space of polynomials with a fixed rational
subspace $V\subset\Q^s$, and the conditions are Taylor coefficients composed with fixed linear maps. The
subspaces $V$ are explicit but have no closed form. Appendix~\ref{app:windows} states this form
(Theorem~\ref{thm:profile}) and describes how its hypotheses are verified.

\subsection{The corank profile for our parameters}\label{ssec:kappa}
The bounds combine into the following statement.

\begin{theorem}\label{thm:profile-D}
For the parameters $(72,11,4,6,4)$ there are $148$ intervals $(a_i,b_i]$ with rational endpoints,
$\frac{185000001}{74000000}=a_1<b_1\le a_2<\dots<b_{148}=\frac{954}{29}$, and a function $\kappa$, affine on each
interval and $0$ outside their union, with the following property. For all large $N$ and every prime $p$ with
$a_iN<p\le b_iN$ for some $i$,
\[
 \nu_N(p)\ge N\kappa(p/N)-2 .
\]
On each interval the bound is given by one of the families of this section, or by the explicit subspaces of
Appendix~\ref{app:windows}. The area of $\kappa$ is the rational number
\begin{multline*}
 \int_0^\infty\kappa(\theta)\,d\theta\\
 =\frac{8933203314101219901240655133431309756867}{258903586359571936582800000000000000000}=\Chi\ldots .
\end{multline*}
\end{theorem}

The first endpoint exceeds $\frac52$ by $\frac{1}{74\cdot10^6}$. Some other endpoints and intercepts are moved by amounts below $10^{-7}$,
to absorb the integer rounding of the families. The intervals are listed in the formalization.

The contributions by ranges of $\theta=p/N$ are as follows. The third column counts the intervals on which
the bound is that of Theorem~\ref{thm:profile} with an explicit subspace.
\begin{center}
\begin{tabular}{@{}lrrl@{}}
\toprule
range of $\theta$ & intervals & with explicit subspace & area \\
\midrule
$(5/2,\,3]$ & $29$ & $9$ & $0.03074$ \\
$(3,\,6]$ & $67$ & $20$ & $0.36842$ \\
$(6,\,12]$ & $31$ & $13$ & $1.06902$ \\
$(12,\,954/29]$ & $21$ & $0$ & $33.03580$ \\
\midrule
total & $148$ & $42$ & $\Chi$ \\
\bottomrule
\end{tabular}
\end{center}
Almost all of the area lies in $(12,\frac{954}{29}]$. There the intervals with $\theta>\frac{74}5$ are as follows.
Adjacent intervals with the same line are merged.
\begin{center}
\begin{tabular}{@{}lll@{}}
\toprule
interval of $\theta$ & $\kappa(\theta)$ & bound \\
\midrule
$(74/5,\,15]$ & $5\theta-74$ & (F2) \\
$(15,\,16]$ & $\theta-14$ & (F2) \\
$(16,\,97/6]$ & $\theta/2-6$ & Theorem~\ref{thm:content}(c), $(p-1)/2-g$ \\
$(97/6,\,182/11]$ & $91-11\theta/2$ & Theorem~\ref{thm:content}(c), second term \\
$(186/11,\,86/5]$ & $11\theta/2-93$ & (F7), law $\mathrm L$ \\
$(86/5,\,87/5]$ & $\theta/2-7$ & (F7), law $\mathrm L$ \\
$(87/5,\,160/9]$ & $80-9\theta/2$ & (F7), law $\mathrm L$ \\
$(132/7,\,97/5]$ & $7\theta/2-66$ & (F7), law $\mathrm H$ with cap \\
$(97/5,\,20]$ & $31-3\theta/2$ & (F7), law $\mathrm H$ \\
$(20,\,222/11]$ & $91-9\theta/2$ & (F7), law $\mathrm H$ \\
$(222/11,\,146/7]$ & $\theta-20$ & (F3) \\
$(146/7,\,43/2]$ & $9\theta/2-93$ & Theorem~\ref{thm:content}(d), $k=4$ \\
$(43/2,\,87/4]$ & $\theta/2-7$ & Theorem~\ref{thm:content}(d), $k=4$ \\
$(87/4,\,160/7]$ & $80-7\theta/2$ & Theorem~\ref{thm:content}(d), $k=4$ \\
$(186/7,\,27]$ & $7\theta/2-93$ & Theorem~\ref{thm:content}(d), $k=3$ \\
$(27,\,28]$ & $5\theta/2-66$ & Theorem~\ref{thm:content}(d), $k=3$ \\
$(28,\,86/3]$ & $3\theta-80$ & Theorem~\ref{thm:content}(a), $3p-(n+q+h)$ \\
$(86/3,\,30]$ & $6$ & Theorem~\ref{thm:content}(a) up to $\theta=29$, then (b) \\
$(30,\,97/3]$ & $66-2\theta$ & Theorem~\ref{thm:content}(b), $n-g-2p$ \\
$(97/3,\,954/29]$ & $66-2\theta$ & (F1) \\
\bottomrule
\end{tabular}
\end{center}
The areas of these rows add up to $33.03580\ldots$, and the $127$ intervals below $12$
contribute $1.46818\ldots$. The exact sum is computed in the
formalization from the list of intervals. Figure~\ref{fig:corank} shows $\kappa$ on $(\frac{74}5,\frac{954}{29}]$.

\begin{figure}[ht]
\centering
\setlength{\unitlength}{1cm}
\begin{picture}(12.4,5.4)(-0.8,-0.9)
\put(0,0){\line(1,0){11.8}}
\put(0,0){\line(0,1){4.9}}
\put(0.000,0){\line(0,-1){0.12}}
\put(0.000,-0.4){\makebox(0,0){$14$}}
\put(1.200,0){\line(0,-1){0.12}}
\put(1.200,-0.4){\makebox(0,0){$16$}}
\put(2.400,0){\line(0,-1){0.12}}
\put(2.400,-0.4){\makebox(0,0){$18$}}
\put(3.600,0){\line(0,-1){0.12}}
\put(3.600,-0.4){\makebox(0,0){$20$}}
\put(4.800,0){\line(0,-1){0.12}}
\put(4.800,-0.4){\makebox(0,0){$22$}}
\put(6.000,0){\line(0,-1){0.12}}
\put(6.000,-0.4){\makebox(0,0){$24$}}
\put(7.200,0){\line(0,-1){0.12}}
\put(7.200,-0.4){\makebox(0,0){$26$}}
\put(8.400,0){\line(0,-1){0.12}}
\put(8.400,-0.4){\makebox(0,0){$28$}}
\put(9.600,0){\line(0,-1){0.12}}
\put(9.600,-0.4){\makebox(0,0){$30$}}
\put(10.800,0){\line(0,-1){0.12}}
\put(10.800,-0.4){\makebox(0,0){$32$}}
\put(0,0.000){\line(-1,0){0.12}}
\put(-0.35,0.000){\makebox(0,0)[r]{$0$}}
\put(0,0.750){\line(-1,0){0.12}}
\put(-0.35,0.750){\makebox(0,0)[r]{$1$}}
\put(0,1.500){\line(-1,0){0.12}}
\put(-0.35,1.500){\makebox(0,0)[r]{$2$}}
\put(0,2.250){\line(-1,0){0.12}}
\put(-0.35,2.250){\makebox(0,0)[r]{$3$}}
\put(0,3.000){\line(-1,0){0.12}}
\put(-0.35,3.000){\makebox(0,0)[r]{$4$}}
\put(0,3.750){\line(-1,0){0.12}}
\put(-0.35,3.750){\makebox(0,0)[r]{$5$}}
\put(0,4.500){\line(-1,0){0.12}}
\put(-0.35,4.500){\makebox(0,0)[r]{$6$}}
\qbezier(0.480,0.000)(0.540,0.375)(0.600,0.750)
\qbezier(0.600,0.750)(0.900,1.125)(1.200,1.500)
\qbezier(1.200,1.500)(1.250,1.531)(1.300,1.562)
\qbezier(1.300,1.562)(1.414,0.781)(1.527,0.000)
\qbezier(1.745,0.000)(1.833,0.600)(1.920,1.200)
\qbezier(1.920,1.200)(1.980,1.237)(2.040,1.275)
\qbezier(2.040,1.275)(2.153,0.637)(2.267,0.000)
\qbezier(2.914,0.000)(3.077,0.712)(3.240,1.425)
\qbezier(3.240,1.425)(3.420,1.087)(3.600,0.750)
\qbezier(3.600,0.750)(3.655,0.443)(3.709,0.136)
\qbezier(3.709,0.136)(3.912,0.390)(4.114,0.643)
\qbezier(4.114,0.643)(4.307,1.728)(4.500,2.812)
\qbezier(4.500,2.812)(4.575,2.859)(4.650,2.906)
\qbezier(4.650,2.906)(4.982,1.453)(5.314,0.000)
\qbezier(7.543,0.000)(7.671,0.562)(7.800,1.125)
\qbezier(7.800,1.125)(8.100,2.062)(8.400,3.000)
\qbezier(8.400,3.000)(8.600,3.750)(8.800,4.500)
\qbezier(8.800,4.500)(8.900,4.500)(9.000,4.500)
\qbezier(9.000,4.500)(9.300,4.500)(9.600,4.500)
\qbezier(9.600,4.500)(10.300,2.750)(11.000,1.000)
\qbezier(11.000,1.000)(11.169,0.578)(11.338,0.155)
\put(11.95,-0.4){\makebox(0,0)[l]{$\theta$}}
\end{picture}

\caption{The corank profile $\kappa(\theta)$ on $(\frac{74}5,\frac{954}{29}]$, where $\nu_N(p)\ge N\kappa(p/N)-2$ for
$p=\theta N$. The bound on the pole count is $f(\theta)=144$ on this whole range, so $\kappa$ is the gain over
OpenAI's denominators.}\label{fig:corank}
\end{figure}

On each interval we use a single bound. The families come from different constructions, and we have not
shown that their kernels are independent, except for the pair (F4). So the bounds are not added at a prime
where several apply.

For Theorem~\ref{thm:content} one checks the hypotheses on the intervals above. The hypothesis $e\ge0$
reads $\theta\ge16$, the hypothesis $n\le4p+2g+1$ reads $\theta\ge15$, and $L\le6p$ reads $\theta\ge29/2$. For
$16\le\theta\le97/6$ one has $\varsigma_0=H+q-6p=101N-6p$, so the first branch of Theorem~\ref{thm:U6}
applies. For larger $\theta$ one has $\varsigma_0=4N$ and the second branch applies, which is positive until
$\theta=182/11$. In Theorem~\ref{thm:S} with $k=4$, $m_1=p$ for $\theta\le43/2$ and $m_1=L-3p=87N-3p$ above, and
$m_0=\max(15N,101N-4p)$, so that $m_0=101N-4p$ for $\theta\le43/2$. The three values $m_1-m_0-e$ are
$9\theta/2-93$, $\theta/2-7$ and $80-7\theta/2$. With $k=3$ one has $m_0=101N-3p$ and
$m_1=\min(p,27N+1)$, and the value $m_1-m_0-e$ is $7\theta/2-93$ on $(\frac{186}7,27]$ and $5\theta/2-66$ on
$(27,28]$, up to the constants $\mp\frac12$. Theorem~\ref{thm:T} gives $3p-(n+q+h)=3\theta-80$ up to the cap
$g=6N$, which it reaches at $\theta=\frac{86}3$. From $\theta=29$ the hypotheses of Theorem~\ref{thm:U}
hold. They read $3p\ge L$, $3p\le H$, $2p\le n+q$ and $2h\le p+1$. The bound is
$\min(n,2p+2g+1)-g-2p-1$, which is $6N$ for $\theta\le30$ and $66N-2p-1$ above. In each case the value of the
theorem is the value in the table minus at most one unit, from the integer rounding. The same kind of check,
with linear inequalities in $N$, $p$ and the parameters, gives the bounds of Section~\ref{ssec:families} on
their intervals, and Appendix~\ref{app:windows} gives the others. This proves Theorem~\ref{thm:profile-D}.
The loss of $2$ in its statement covers the integer rounding of all the bounds.
\section{The neighbourhood of $G$}\label{sec:real}

This section and the next two prove that $\Delta_N$ is small on a whole neighbourhood of $G$. The idea is to
write the perturbation $\Delta_N(x)-\Delta_N(G)$ as a sum of determinants in which some columns of $\cA_N$
are replaced by values of the row functions at prescribed nodes. Each such determinant is bounded by
OpenAI's energy method, extended to configurations with prescribed points (Section~\ref{sec:energy}). The
loss caused by each replaced column, in the exponent, determines $\rho$. OpenAI's method has two cases. In
the first it already gives the improved decay rate. In the second it falls short when few columns are
replaced, and Section~\ref{sec:signed} closes that gap.

\subsection{The statement}
OpenAI's energy method bounds $|\Delta_N(G)|$ by $\exp\bigl(-(n_0^2(U-\frac12\log2)+o(1))N^2\bigr)$, where $U$
is an upper bound for a discrete logarithmic energy (Section~\ref{sec:energy}). For our parameters we
certify $U=\Utarget$, so OpenAI's estimate gives the decay rate $D_0=n_0^2(U-\frac12\log2)=\DecayOA\ldots$
Let
\begin{gather*}
 \Gamma_0=\Gain,\qquad U_0=U+\Gamma_0/n_0^2=\Uzero\ldots,\\
 D=n_0^2(U_0-\tfrac12\log2)=\Decay\ldots
\end{gather*}
The number $\Gamma_0$ is the improvement of the decay rate proved in Section~\ref{sec:signed}. Put
\begin{equation}\label{eq:rho-def}
 \rho=n_0c_2+W_0,\qquad c_2=\Ctwo,\quad W_0=\WZero,\quad \rho=\Rho .
\end{equation}
Here $c_2$ is the increase of the energy bound per prescribed column (Section~\ref{sec:energy}), and $W_0$
bounds the logarithmic size of the weight and the binomial factor of one column, per unit of $N$
(\S\ref{ssec:weights}).

\begin{theorem}\label{thm:real}
For our parameters, uniformly for real $x$ with $|x-G|\le e^{-\rho N}$,
\[
 |\Delta_N(x)|\le e^{-(D+o(1))N^2}\qquad(N\to\infty).
\]
\end{theorem}

We use the following abbreviations for general parameters:
\begin{gather*}
 \alpha_A=\frac An,\quad \beta_h=\frac{2h_0}{n_0},\quad \gamma_H=\frac Hn,\quad \gamma_d=\frac{d_0}{n_0},\\
 a_0=1-\frac{b_0+q_0+g_0}{n_0},\quad a_1=3-\frac{b_0+q_0+2g_0}{n_0}.
\end{gather*}
For our parameters these are $\alpha_A=11/29$, $\beta_h=3/29$, $\gamma_H=39/29$, $\gamma_d=53/58$,
$a_0=41/58$ and $a_1=76/29$.

\subsection{A quadrature formula for the coefficient of $G$}\label{ssec:quadrature}
Let $\varepsilon=x-G$. Then $\Delta_N(x)=\det(\cA_N+\varepsilon\cL_N)$. The question is how large the columns
of $\cL_N$ are compared with those of $\cA_N$. By \eqref{eq:L} the entries of $\cL_N$ are linear
functionals of the coefficients of the row polynomials. These polynomials have bounded degree, so the
functionals can be written exactly as finite combinations of point evaluations. Let $t_1,\ldots,t_J$ be
distinct points, with $J$ at least the number of degrees of freedom of the rows. Lagrange interpolation
gives
\begin{equation}\label{eq:quad}
 \cL_N(r,k)=\sum_{a=1}^JA_r(t_a)\,\beta_{a,k},
\end{equation}
where the weight $\beta_{a,k}$ is the value of the same functional on the $a$-th Lagrange basis polynomial
$\ell_a$. This is a quadrature formula of Gauss type. It makes the columns of $\cL_N$ combinations of point
values of the row functions that define $\cA_N$. So the perturbed determinant can be treated by the energy
method. The price is the size of the weights.

We take $J=d_0N=53N$ nodes. Put $g_*=401/2500$, $a_*=2g_*/(1+g_*^2)=2005000/6410801$ and $J'=\lceil J/2\rceil$.
The nodes are the points
\[
 t_{j,\pm}=\pm\sqrt{a_*^2+(1-a_*^2)(1+s_j)/2},\qquad T_{J'}(s_j)=0,
\]
that is, the roots of $T_{J'}\bigl((2t^2-(1+a_*^2))/(1-a_*^2)\bigr)$. We keep all $J'$ positive roots and the
first $J-J'$ negative ones, so for odd $J$ exactly one root is dropped. The nodes lie in
$S=(-1,-a_*)\cup(a_*,1)$, and in the coordinate $x$ with $t=2x/(1+x^2)$ every node has $|x|>g_*$. The nodes
avoid a neighbourhood of $x=0$, where the fields of the energy method are largest, and they are Chebyshev
nodes for $S$. Dropping a root changes the nodal polynomial by a linear factor, so it changes the
derivative bound and the size of the Lagrange polynomials only by factors polynomial in $N$. So an odd
number of nodes does not change the exponent. This is needed, since $J=53N$ is odd for odd $N$.

Let $g_S$ be the Green function of $\C\setminus\bar S$ with pole at infinity. In closed form
\[
 g_S(z)=\tfrac12\operatorname{arccosh}\Bigl(\frac{|z^2-1|+|z^2-a_*^2|}{1-a_*^2}\Bigr).
\]
The Lagrange polynomials satisfy $|\ell_a(z)|\le e^{53Ng_S(z)+o(N)}$ at distance at least $c/N$ from the
nodes, by the derivative identity for Chebyshev polynomials at their roots. The weight of column $k$ is a
coefficient functional of $\ell_a$ multiplied by the column function. After removal of the common factor of
the rows it equals
\[
 \beta'_k=4\,[t^0]\,F_k(t)\,(1-t^{-2})^{-1/2},\qquad F_k(t)=t^{\,A-b-q-k}(1-t)^h(t-1)^q\ell_a(t),
\]
where $(1-t^{-2})^{-1/2}=\sum_lc_lt^{-2l}$ with $c_l=\binom{2l}l4^{-l}$. The next lemma bounds it by a contour
that depends on $k$.

\begin{lemma}[keyhole bound]\label{lem:keyhole}
Let $P$ be a real polynomial, $E\ge0$, and $F(t)=t^{-E}P(t)$. Let $0<r\le1$, and suppose that $|F|\le M$ on
the circle $|t|=r$ and $|F|\le M'$ on the segments $r\le|t|\le1$ of the real axis. Then
\[
 \Bigl|4\sum_{l\ge0}c_l\,[t^{E+2l}]P\Bigr|\le4(1+E)\max(M,M').
\]
\end{lemma}

\begin{proof}
Write $F=\Pi+R$ with $\Pi$ a polynomial and $R=\sum_{i=1}^E\pi_it^{-i}$ the principal part at $0$. The
functional sees only $\Pi$. Since $c_l=\frac1\pi\int_{-1}^1t^{2l}(1-t^2)^{-1/2}dt$, its value is
$\frac4\pi\int_{-1}^1\Pi(t)(1-t^2)^{-1/2}dt$, which is at most $4\max_{[-1,1]}|\Pi|$. By Cauchy's estimate
$|\pi_i|\le Mr^i$, so $|R(t)|\le EM$ for $r\le|t|\le1$, and $|\Pi|\le M'+EM$ there. For $|t|<r$ the maximum
principle gives $|\Pi(t)|\le\max_{|t'|=r}|\Pi|\le(1+E)M$.
\end{proof}

The formal proof replaces the integral by the finite Gauss--Chebyshev identity
$\frac1M\sum_{j<M}\cos^{2l}\frac{(2j+1)\pi}{2M}=c_l$ for $l<M$, and Cauchy's estimate by the finite Cauchy formula
at roots of unity. So no integration is needed. For $0\le y\le1$ and $z\in\C$ let
\[
 \Phi_y(z)=(10-58y)\log|z|+6\log|1-z|+53g_S(z).
\]
Here $10=(A-b-q)/N$ and $6=(q+h)/N$. With $y=k/n$, Lemma~\ref{lem:keyhole} gives
\begin{equation}\label{eq:wk}
 |\beta'_k|\le e^{(w(y)+o(1))N},\qquad w(y)=\min_r\max_{z\in\Gamma_r}\Phi_y(z),
\end{equation}
where $\Gamma_r$ is the circle $|z|=r$ together with the segments $r\le|z|\le1$ of the real axis. For $y$ near
$1$ the factor $|z|^{10-58y}$ favours a large circle, and for small $y$ a small one. The circle $|z|=1.001$ alone
gives $w(y)\le52.79$ for every $y$. The keyhole contour improves this only slightly at the worst column. Its
use is elsewhere: it lowers $w(y)$ at the columns where the binomial factor of Lemma~\ref{lem:binomvand} is
largest.

\subsection{Column expansion and Andr\'eief's identity}\label{ssec:andreief}
Write $\mathcal P(r,a)=A_r(t_a)$. Expanding $\det(\cA_N+\varepsilon\mathcal P\beta)$ multilinearly in the
columns gives a sum over maps $f$ from the columns to $\{\varnothing,1,\ldots,J\}$ of
$\prod_{f(k)\ne\varnothing}\varepsilon\beta_{f(k),k}\det M_f$. Here $M_f$ is the matrix whose $k$-th column is
the $k$-th column of $\cA_N$ if $f(k)=\varnothing$, and the vector $(A_r(t_{f(k)}))_r$ otherwise.

\begin{lemma}\label{lem:expansion}
Suppose that $|\det M_f|\le M_0\prod_{f(k)\ne\varnothing}C_k$ for every $f$, and that $|\beta_{a,k}|\le B_k$.
Then
\[
 \bigl|\det(\cA_N+\varepsilon\mathcal P\beta)-\det\cA_N\bigr|\le M_0\Bigl(\prod_k(1+J|\varepsilon|B_kC_k)-1\Bigr).
\]
\end{lemma}

The lemma allows a different constant for each column. Each $M_f$ has $n'=n-m$ columns that are integrals
and $m$ columns that are point values. Andr\'eief's identity extends to this situation.

\begin{lemma}\label{lem:andreief}
Let $m_1,m_2$ be finite measures, and let $F_r$ $(r<n'+m)$, $F'_i$ $(i<n')$ and $k$ be integrable functions.
Let $B$ be the matrix with $B(r,i)=\iint F_r(t)k(t,s)F'_i(s)\,dm_1(t)\,dm_2(s)$ for $i<n'$ and
$B(r,n'+j)=F_r(t'_j)$ for $j<m$. Then
\begin{multline*}
 (n'!)^2\det B=\int\!\!\int\det\bigl[F_r((t\Vert t')_c)\bigr]_{r,c}\\
 \cdot\det[k(t_i,s_j)]_{i,j<n'}\det[F'_i(s_j)]_{j,i<n'}\,dm_1^{n'}(t)\,dm_2^{n'}(s),
\end{multline*}
where $t\Vert t'$ is the list of the $n'$ variables followed by the $m$ prescribed points.
\end{lemma}

\begin{proof}
For $m=0$ this is Andr\'eief's identity applied twice \cite{Forrester2018Andreief}. In general, expand
$\det B$ multilinearly in its first $n'$ columns. This introduces one pair of variables $(t_i,s_i)$ for each
of them. Then antisymmetrize, first in the $t$ variables and then in the $s$ variables. The last $m$ columns
are constant throughout.
\end{proof}

With $k(t,s)=1/(1-ts)$ the middle determinant is a Cauchy determinant. The third is a generalized
Vandermonde determinant $\prod_js_j^b(1-s_j)^q\det[s_j^{E_i}]$. Here the exponents $E_i<n$ are those of the
columns that remain integrals, and the missing exponents are the replaced columns. We bound the generalized
Vandermonde determinant by a product over the missing exponents.

\begin{lemma}\label{lem:binomvand}
Let $s_1,\ldots,s_{n'}\in[0,1]$, let $E_1<\dots<E_{n'}<n$ be integers, and let $\{k_1,\ldots,k_m\}$ be the
complement of $\{E_i\}$ in $\{0,\ldots,n-1\}$. Then
\[
 \bigl|\det[s_l^{E_i}]\bigr|\le\prod_{j=1}^m\mathcal B(k_j)\prod_{i<l}|s_i-s_l|,\qquad
 \mathcal B(k)=\sum_{i=0}^{m-1}\binom{n'}{n'-k+i},
\]
with $\binom{n'}r=0$ for $r\notin[0,n']$.
\end{lemma}

\begin{proof}
Let $P(X)=\prod_l(X-s_l)=\sum_r(-1)^re_r(s)X^{n'-r}$. Multiplying the matrix by a unitriangular matrix built
from $P$ gives $|\det[s_l^{E_i}]|=\prod_{i<l}|s_i-s_l|\cdot|\det Q|$, where $Q$ is the $m\times m$ matrix with
entries $Q_{j,i}=[X^{k_j-i}]P$. So $|Q_{j,i}|=e_{n'-k_j+i}(s)\le\binom{n'}{n'-k_j+i}$, since $s\in[0,1]^{n'}$.
Hadamard's inequality, applied to the rows of $Q$ in the $\ell^1$ norm, gives the claim.
\end{proof}

\subsection{The weight of one column}\label{ssec:weights}
The bound of Lemma~\ref{lem:binomvand} does not depend on the points $s_l$. Since $\binom{n'}r\le\binom nr$,
\[
 \log\mathcal B(k)\le n_0\tilde h_2(k/n)N+O(\log N),
\]
where $\tilde h_2(y)$ is the binary entropy $h_2(y)=-y\log y-(1-y)\log(1-y)$ for $y\ge\frac12$ and $\log2$ for
$y<\frac12$. So the replaced column $k$ contributes at most $e^{(w(y)+n_0\tilde h_2(y)+o(1))N}$, with $y=k/n$.
Put
\[
 W_0\ge\sup_{0\le y\le1}\bigl(w(y)+58\tilde h_2(y)\bigr).
\]
Appendix~\ref{app:field} proves $W_0=\WZero$.

The two factors peak at different columns. The weight $w(y)$ increases up to $y\approx0.84$, while
$\tilde h_2(y)$ is largest on $[0,\frac12]$. So the maximum of $w+58\tilde h_2$ is about $87.96$, at
$y\approx0.62$, well below the sum $52.79+40.20$ of the separate maxima (Figure~\ref{fig:columns}). This saving of
about $5$ in $\rho$ is the reason for treating the columns one at a time.

\begin{figure}[ht]
\centering
\setlength{\unitlength}{1cm}
\begin{picture}(12.6,5.6)(-1.0,-0.8)
\put(0,0){\line(1,0){10.2}}
\put(0,0){\line(0,1){4.9}}
\put(0.000,0){\line(0,-1){0.12}}
\put(0.000,-0.4){\makebox(0,0){$0$}}
\put(2.000,0){\line(0,-1){0.12}}
\put(2.000,-0.4){\makebox(0,0){$0.2$}}
\put(4.000,0){\line(0,-1){0.12}}
\put(4.000,-0.4){\makebox(0,0){$0.4$}}
\put(6.000,0){\line(0,-1){0.12}}
\put(6.000,-0.4){\makebox(0,0){$0.6$}}
\put(8.000,0){\line(0,-1){0.12}}
\put(8.000,-0.4){\makebox(0,0){$0.8$}}
\put(10.000,0){\line(0,-1){0.12}}
\put(10.000,-0.4){\makebox(0,0){$1$}}
\put(0,0.000){\line(-1,0){0.12}}
\put(-0.25,0.000){\makebox(0,0)[r]{$0$}}
\put(0,1.000){\line(-1,0){0.12}}
\put(-0.25,1.000){\makebox(0,0)[r]{$20$}}
\put(0,2.000){\line(-1,0){0.12}}
\put(-0.25,2.000){\makebox(0,0)[r]{$40$}}
\put(0,3.000){\line(-1,0){0.12}}
\put(-0.25,3.000){\makebox(0,0)[r]{$60$}}
\put(0,4.000){\line(-1,0){0.12}}
\put(-0.25,4.000){\makebox(0,0)[r]{$80$}}
\put(10.15,-0.4){\makebox(0,0)[l]{$y$}}
\qbezier(0.000,0.586)(0.100,0.621)(0.200,0.656)
\qbezier(0.200,0.656)(0.300,0.691)(0.400,0.725)
\qbezier(0.400,0.725)(0.500,0.760)(0.600,0.795)
\qbezier(0.600,0.795)(0.700,0.830)(0.800,0.865)
\qbezier(0.800,0.865)(0.900,0.900)(1.000,0.935)
\qbezier(1.000,0.935)(1.100,0.970)(1.200,1.005)
\qbezier(1.200,1.005)(1.300,1.040)(1.400,1.075)
\qbezier(1.400,1.075)(1.500,1.110)(1.600,1.144)
\qbezier(1.600,1.144)(1.700,1.179)(1.800,1.214)
\qbezier(1.800,1.214)(1.900,1.249)(2.000,1.284)
\qbezier(2.000,1.284)(2.100,1.319)(2.200,1.354)
\qbezier(2.200,1.354)(2.300,1.389)(2.400,1.424)
\qbezier(2.400,1.424)(2.500,1.459)(2.600,1.494)
\qbezier(2.600,1.494)(2.700,1.529)(2.800,1.563)
\qbezier(2.800,1.563)(2.900,1.598)(3.000,1.633)
\qbezier(3.000,1.633)(3.100,1.668)(3.200,1.703)
\qbezier(3.200,1.703)(3.300,1.738)(3.400,1.773)
\qbezier(3.400,1.773)(3.500,1.808)(3.600,1.843)
\qbezier(3.600,1.843)(3.700,1.877)(3.800,1.911)
\qbezier(3.800,1.911)(3.900,1.944)(4.000,1.976)
\qbezier(4.000,1.976)(4.100,2.006)(4.200,2.037)
\qbezier(4.200,2.037)(4.300,2.065)(4.400,2.094)
\qbezier(4.400,2.094)(4.500,2.121)(4.600,2.148)
\qbezier(4.600,2.148)(4.700,2.174)(4.800,2.199)
\qbezier(4.800,2.199)(4.900,2.223)(5.000,2.247)
\qbezier(5.000,2.247)(5.100,2.269)(5.200,2.291)
\qbezier(5.200,2.291)(5.300,2.312)(5.400,2.333)
\qbezier(5.400,2.333)(5.500,2.353)(5.600,2.372)
\qbezier(5.600,2.372)(5.700,2.390)(5.800,2.408)
\qbezier(5.800,2.408)(5.900,2.425)(6.000,2.442)
\qbezier(6.000,2.442)(6.100,2.457)(6.200,2.472)
\qbezier(6.200,2.472)(6.300,2.486)(6.400,2.501)
\qbezier(6.400,2.501)(6.500,2.513)(6.600,2.526)
\qbezier(6.600,2.526)(6.700,2.537)(6.800,2.549)
\qbezier(6.800,2.549)(6.900,2.559)(7.000,2.569)
\qbezier(7.000,2.569)(7.100,2.578)(7.200,2.587)
\qbezier(7.200,2.587)(7.300,2.594)(7.400,2.602)
\qbezier(7.400,2.602)(7.500,2.608)(7.600,2.614)
\qbezier(7.600,2.614)(7.700,2.619)(7.800,2.624)
\qbezier(7.800,2.624)(7.900,2.628)(8.000,2.631)
\qbezier(8.000,2.631)(8.100,2.633)(8.200,2.636)
\qbezier(8.200,2.636)(8.300,2.636)(8.400,2.637)
\qbezier(8.400,2.637)(8.500,2.637)(8.600,2.637)
\qbezier(8.600,2.637)(8.700,2.637)(8.800,2.637)
\qbezier(8.800,2.637)(8.900,2.637)(9.000,2.637)
\qbezier(9.000,2.637)(9.100,2.637)(9.200,2.637)
\qbezier(9.200,2.637)(9.300,2.637)(9.400,2.637)
\qbezier(9.400,2.637)(9.500,2.637)(9.600,2.637)
\qbezier(9.600,2.637)(9.700,2.637)(9.800,2.637)
\qbezier(9.800,2.637)(9.900,2.637)(10.000,2.637)
\qbezier(0.000,2.010)(0.100,2.010)(0.200,2.010)
\qbezier(0.200,2.010)(0.300,2.010)(0.400,2.010)
\qbezier(0.400,2.010)(0.500,2.010)(0.600,2.010)
\qbezier(0.600,2.010)(0.700,2.010)(0.800,2.010)
\qbezier(0.800,2.010)(0.900,2.010)(1.000,2.010)
\qbezier(1.000,2.010)(1.100,2.010)(1.200,2.010)
\qbezier(1.200,2.010)(1.300,2.010)(1.400,2.010)
\qbezier(1.400,2.010)(1.500,2.010)(1.600,2.010)
\qbezier(1.600,2.010)(1.700,2.010)(1.800,2.010)
\qbezier(1.800,2.010)(1.900,2.010)(2.000,2.010)
\qbezier(2.000,2.010)(2.100,2.010)(2.200,2.010)
\qbezier(2.200,2.010)(2.300,2.010)(2.400,2.010)
\qbezier(2.400,2.010)(2.500,2.010)(2.600,2.010)
\qbezier(2.600,2.010)(2.700,2.010)(2.800,2.010)
\qbezier(2.800,2.010)(2.900,2.010)(3.000,2.010)
\qbezier(3.000,2.010)(3.100,2.010)(3.200,2.010)
\qbezier(3.200,2.010)(3.300,2.010)(3.400,2.010)
\qbezier(3.400,2.010)(3.500,2.010)(3.600,2.010)
\qbezier(3.600,2.010)(3.700,2.010)(3.800,2.010)
\qbezier(3.800,2.010)(3.900,2.010)(4.000,2.010)
\qbezier(4.000,2.010)(4.100,2.010)(4.200,2.010)
\qbezier(4.200,2.010)(4.300,2.010)(4.400,2.010)
\qbezier(4.400,2.010)(4.500,2.010)(4.600,2.010)
\qbezier(4.600,2.010)(4.700,2.010)(4.800,2.010)
\qbezier(4.800,2.010)(4.900,2.010)(5.000,2.010)
\qbezier(5.000,2.010)(5.100,2.009)(5.200,2.008)
\qbezier(5.200,2.008)(5.300,2.004)(5.400,2.001)
\qbezier(5.400,2.001)(5.500,1.995)(5.600,1.989)
\qbezier(5.600,1.989)(5.700,1.981)(5.800,1.973)
\qbezier(5.800,1.973)(5.900,1.962)(6.000,1.952)
\qbezier(6.000,1.952)(6.100,1.939)(6.200,1.926)
\qbezier(6.200,1.926)(6.300,1.910)(6.400,1.895)
\qbezier(6.400,1.895)(6.500,1.877)(6.600,1.859)
\qbezier(6.600,1.859)(6.700,1.838)(6.800,1.818)
\qbezier(6.800,1.818)(6.900,1.795)(7.000,1.772)
\qbezier(7.000,1.772)(7.100,1.746)(7.200,1.720)
\qbezier(7.200,1.720)(7.300,1.691)(7.400,1.662)
\qbezier(7.400,1.662)(7.500,1.630)(7.600,1.598)
\qbezier(7.600,1.598)(7.700,1.563)(7.800,1.528)
\qbezier(7.800,1.528)(7.900,1.490)(8.000,1.451)
\qbezier(8.000,1.451)(8.100,1.409)(8.200,1.367)
\qbezier(8.200,1.367)(8.300,1.321)(8.400,1.275)
\qbezier(8.400,1.275)(8.500,1.225)(8.600,1.174)
\qbezier(8.600,1.174)(8.700,1.119)(8.800,1.064)
\qbezier(8.800,1.064)(8.900,1.003)(9.000,0.943)
\qbezier(9.000,0.943)(9.100,0.876)(9.200,0.808)
\qbezier(9.200,0.808)(9.300,0.733)(9.400,0.658)
\qbezier(9.400,0.658)(9.500,0.573)(9.600,0.487)
\qbezier(9.600,0.487)(9.700,0.386)(9.800,0.284)
\qbezier(9.800,0.284)(9.900,0.142)(10.000,0.000)
\qbezier(0.000,2.596)(0.100,2.631)(0.200,2.666)
\qbezier(0.200,2.666)(0.300,2.701)(0.400,2.736)
\qbezier(0.400,2.736)(0.500,2.771)(0.600,2.805)
\qbezier(0.600,2.805)(0.700,2.840)(0.800,2.875)
\qbezier(0.800,2.875)(0.900,2.910)(1.000,2.945)
\qbezier(1.000,2.945)(1.100,2.980)(1.200,3.015)
\qbezier(1.200,3.015)(1.300,3.050)(1.400,3.085)
\qbezier(1.400,3.085)(1.500,3.120)(1.600,3.155)
\qbezier(1.600,3.155)(1.700,3.190)(1.800,3.224)
\qbezier(1.800,3.224)(1.900,3.259)(2.000,3.294)
\qbezier(2.000,3.294)(2.100,3.329)(2.200,3.364)
\qbezier(2.200,3.364)(2.300,3.399)(2.400,3.434)
\qbezier(2.400,3.434)(2.500,3.469)(2.600,3.504)
\qbezier(2.600,3.504)(2.700,3.539)(2.800,3.574)
\qbezier(2.800,3.574)(2.900,3.608)(3.000,3.643)
\qbezier(3.000,3.643)(3.100,3.678)(3.200,3.713)
\qbezier(3.200,3.713)(3.300,3.748)(3.400,3.783)
\qbezier(3.400,3.783)(3.500,3.818)(3.600,3.853)
\qbezier(3.600,3.853)(3.700,3.887)(3.800,3.921)
\qbezier(3.800,3.921)(3.900,3.954)(4.000,3.986)
\qbezier(4.000,3.986)(4.100,4.016)(4.200,4.047)
\qbezier(4.200,4.047)(4.300,4.076)(4.400,4.104)
\qbezier(4.400,4.104)(4.500,4.131)(4.600,4.158)
\qbezier(4.600,4.158)(4.700,4.184)(4.800,4.209)
\qbezier(4.800,4.209)(4.900,4.233)(5.000,4.257)
\qbezier(5.000,4.257)(5.100,4.278)(5.200,4.299)
\qbezier(5.200,4.299)(5.300,4.317)(5.400,4.334)
\qbezier(5.400,4.334)(5.500,4.348)(5.600,4.361)
\qbezier(5.600,4.361)(5.700,4.371)(5.800,4.381)
\qbezier(5.800,4.381)(5.900,4.387)(6.000,4.393)
\qbezier(6.000,4.393)(6.100,4.396)(6.200,4.398)
\qbezier(6.200,4.398)(6.300,4.397)(6.400,4.395)
\qbezier(6.400,4.395)(6.500,4.390)(6.600,4.385)
\qbezier(6.600,4.385)(6.700,4.376)(6.800,4.367)
\qbezier(6.800,4.367)(6.900,4.354)(7.000,4.341)
\qbezier(7.000,4.341)(7.100,4.323)(7.200,4.306)
\qbezier(7.200,4.306)(7.300,4.285)(7.400,4.264)
\qbezier(7.400,4.264)(7.500,4.238)(7.600,4.212)
\qbezier(7.600,4.212)(7.700,4.182)(7.800,4.152)
\qbezier(7.800,4.152)(7.900,4.117)(8.000,4.082)
\qbezier(8.000,4.082)(8.100,4.042)(8.200,4.003)
\qbezier(8.200,4.003)(8.300,3.957)(8.400,3.912)
\qbezier(8.400,3.912)(8.500,3.862)(8.600,3.812)
\qbezier(8.600,3.812)(8.700,3.756)(8.800,3.701)
\qbezier(8.800,3.701)(8.900,3.641)(9.000,3.580)
\qbezier(9.000,3.580)(9.100,3.513)(9.200,3.446)
\qbezier(9.200,3.446)(9.300,3.370)(9.400,3.295)
\qbezier(9.400,3.295)(9.500,3.210)(9.600,3.124)
\qbezier(9.600,3.124)(9.700,3.023)(9.800,2.921)
\qbezier(9.800,2.921)(9.900,2.779)(10.000,2.637)
\multiput(0,4.650)(0.2,0){51}{\line(1,0){0.08}}
\put(0.2,4.870){\makebox(0,0)[l]{$52.79+40.20$}}
\put(6.200,4.398){\makebox(0,0){$\bullet$}}
\put(6.200,4.048){\makebox(0,0){$87.96$}}
\put(10.15,2.637){\makebox(0,0)[l]{$w(y)$}}
\put(2.0,2.260){\makebox(0,0)[l]{$58\tilde h_2(y)$}}
\put(2.6,3.804){\makebox(0,0)[r]{sum}}
\end{picture}

\caption{The weight $w(y)$ of column $k=yn$, the binomial factor $58\tilde h_2(y)$, and their sum, per unit
of $N$. The values were computed numerically on a grid of step $0.02$ and are not part of the proof. The
certified bound for the sum is $W_0=\WZero$. The dotted line is the sum of the separate maxima.}
\label{fig:columns}
\end{figure}

The remaining factor of $\rho$ is $n_0c_2$, the increase of the energy bound caused by the prescribed
columns. It comes from Section~\ref{sec:energy}, and with it $\rho=58\cdot\Ctwo+\WZero=\Rho$.

\section{The energy method with prescribed columns}\label{sec:energy}

This section recalls OpenAI's energy method, extends it to determinants with prescribed columns, and
proves Theorem~\ref{thm:real} assuming the bound of Section~\ref{sec:signed}. Throughout, $\alpha_1=\AlphaOne$.

\subsection{OpenAI's energy}\label{ssec:energy}
By Andr\'eief's identity, $\det\cA_N$ is an integral over $n$ points $x_i\in(-1,1)$ and $n$ points
$s_j\in(0,1)$. OpenAI bounds the integrand by $e^{n^2E}$, where $E$ is a discrete logarithmic energy of the
points, written in the coordinate $x$ with $t=2x/(1+x^2)$ \cite[\S5]{OpenAI2026}. It has the form
\begin{multline*}
 n^2E=\sum_iV_1(x_i)+\sum_jV_2(s_j)+\sum_{i<i'}\bigl(\log|x_i-x_{i'}|+\log(1-x_ix_{i'})\bigr)\\
 +\sum_{j<j'}\log|s_j-s_{j'}|-\sum_{i,j}\log(1-2x_is_j+x_i^2),
\end{multline*}
schematically, with external fields $V_1,V_2$ of order $n$. The first two sums are the
external fields. The next two are the logarithmic interactions among points of the same kind, and the last
is the interaction between the two kinds. There are two cases, according to which of OpenAI's two bounds for
the determinant of the row functions is used \cite[Props.~5.2, 5.3]{OpenAI2026}. The case with doubled interaction is the one in which OpenAI's condition
$\sum_i(1-x_i^2)/(1+x_i^2)\le n-1-2g$ \cite[Eq.~(67)]{OpenAI2026} holds, and the constrained case is its complement. In the \emph{constrained
case} the row field acquires a term with a multiplier $\lambda_1$ for the constraint that defines the case.
In the \emph{case with doubled interaction} the interaction among the $x_i$ has weight $2$ instead of $1$. This case also bounds a determinant
$\det\mathsf L$ of size $n$ by $e^{O(n\log n)}$, which is the factor that Section~\ref{sec:signed} improves.

The energy is bounded from above by duality. Given trial sequences $u=(u_k)$ and $v=(v_k)$, one writes the
interaction terms as quadratic forms in the moments $\frac1n\sum_iT_k(x_i)$ and $\frac1n\sum_js_j^k$. Then one
completes the squares against $u$ and $v$. What remains is the sum of the norms $\|u\|^2=\sum_ku_k^2/k$ and
$\|v\|^2$ and, for each point, a function of one variable. With $\mathsf T(u,x)=\sum_ku_kT_k(x)/k$ and
$\mathsf S(v,x)=\sum_kv_kx^k/k$, these functions are, in the case with doubled interaction,
\begin{align*}
 X(x)&=\alpha_A\log|x|+\beta_h\log(1-x)-\gamma_H\log(1+x^2)-4\mathsf T(u,x)-\mathsf S(v,x),\\
 Y(s)&=\tfrac{b_0}{n_0}\log s+\tfrac{q_0}{n_0}\log(1-s)+2\mathsf T(v,s).
\end{align*}
In the constrained case, with a trial sequence $u'$ in place of $u$, they are
\begin{multline*}
 X_1(x)=\alpha_A\log|x|+\beta_h\log(1-x)-(\gamma_H-\tfrac12)\log(1+x^2)\\
 +\lambda_1\bigl(\tfrac{2g_0}{n_0}-\tfrac{2x^2}{1+x^2}\bigr)-2\mathsf T(u',x)-\mathsf S(v,x),
\end{multline*}
with the same $Y$. If $X\le M_X$ and $Y\le M_Y$ on the whole domains, the energy of every configuration is at
most a constant determined by $M_X$, $M_Y$ and the norms. In the language of potential theory the extremal
configuration is described by an equilibrium problem for two interacting measures in an external field
\cite{SaffTotik}. The best trial sequences are the moments of its solution, as Carreira observed
\cite{Carreira2026}. Any trial sequences give a valid upper bound.

\subsection{Prescribed points}\label{ssec:prescribed}
When $m$ of the $n$ columns are point values, only $n-m$ points of each kind are integrated over, and the
energy is still normalized by $n$. The prescribed points enter through the row determinant, now evaluated
at the free points together with the prescribed ones. There are two kinds of prescribed points: the
quadrature nodes, which lie in the set $E_x$ of Section~\ref{sec:real}, and points $y\in(0,1)$. Let $\alpha=m/n$ be the fraction of
prescribed columns and $\alpha_2$ the fraction of prescribed points of the second kind. Each missing free point
changes the energy bound by a bounded amount. The total change is linear in $m$, and its slope $c_2$ enters
$\rho$.

Every step of OpenAI's argument remains valid uniformly in the number of free points. The regularization of
the logarithmic interactions, the diagonal corrections and the completion of squares only involve sums over
the points present. For the prescribed points of the first kind the relevant field is
\[
 F_1(x)=-\gamma_d\log(1+x^2)-2\mathsf T(u,x),\qquad x\in E_x .
\]
For those of the second kind it is
\[
 F_2(y)=-\gamma_d\log(1+y^2)+\frac{n_0-2g_0-1}{n_0}\log y-u_1y,\qquad0<y<1,
\]
where $u_1$ is the first coefficient of $u$. The coefficient $(n_0-2g_0-1)/n_0$ is a lower bound, valid for all
$N\ge1$, of $(n-1-2g)/n$. By Lemma~\ref{lem:binomvand} the generalized Vandermonde determinant contributes only
the factor $\prod_{i<l}|s_i-s_l|$ to the energy, and the binomial factors are assigned to the columns. The
energy bound then reduces to a quadratic inequality in $\alpha$ and $\alpha_2$.

\begin{lemma}[energy inequality with prescribed points]\label{lem:energy-ineq}
Let $\ell=\log2$, $\alpha_1>0$ and $u\ge0$, and let $\mathcal N,\mathcal F,\mathcal R,\mathcal F_2$ and $c_2$ be real
numbers with
\begin{gather*}
 \mathcal N+\mathcal F+U_0-a_0\ell\le u,\qquad \ell-\mathcal R+\mathcal F_2\le0,\\
 (a_1-1)\ell-\mathcal F+\mathcal R-c_2-\ell\alpha_1+\frac u{\alpha_1}\le0 .
\end{gather*}
Then for all $\alpha\ge\alpha_1$ and $\alpha_2\ge0$
\begin{multline*}
 \bigl(-(a_0-\tfrac12)+(a_1-1)\alpha-\alpha^2+\alpha_2-\alpha_2^2\bigr)\ell+\mathcal N+(1-\alpha)\mathcal F\\
 +(\alpha-\alpha_2)\mathcal R+\alpha_2\mathcal F_2\le-U_0+\tfrac12\ell+c_2\alpha .
\end{multline*}
\end{lemma}

\begin{proof}
Put $s=(a_1-1)\ell-\mathcal F+\mathcal R-c_2$. The left side minus the right side equals
\[
 (\mathcal N+\mathcal F+U_0-a_0\ell)+\alpha s-\ell\alpha^2+\alpha_2(\ell-\mathcal R+\mathcal F_2)-\ell\alpha_2^2 .
\]
The last two terms are nonpositive by the second hypothesis. The first three are at most
$\alpha(u/\alpha+s-\ell\alpha)$. The function $u/\alpha-\ell\alpha$ decreases in $\alpha$, so this is at most
$\alpha(u/\alpha_1+s-\ell\alpha_1)\le0$.
\end{proof}

The lemma is applied with $\mathcal N$ a bound for $2\|u\|^2+\frac12\|v\|^2$, $\mathcal F=M_X+M_Y$,
$\mathcal R$ a bound for the supremum of $F_1$ and $\mathcal F_2$ a bound for the supremum of $F_2$. The
analogue for the constrained case has no prescribed points of the second kind. Its polynomial in $\alpha$ is
\[
 \bigl(-(a_0-\tfrac12)+(a_1-1)\alpha-\alpha^2\bigr)\ell+n_1+(1-\alpha)\mathcal F_1+\alpha\lambda_1\tfrac{2g_0}{n_0}
 \le-U_0+\tfrac12\ell+c_2\alpha,
\]
with $n_1$ a bound for the norms of the trial sequences and $\mathcal F_1=M_{X_1}+M_Y$. The coefficient of $\alpha^2$ is negative, so this holds for all real $\alpha$ as soon as the discriminant of the
difference of the two sides is nonpositive.

For our parameters the trial sequences of Appendix~\ref{app:field} give the hypotheses of
Lemma~\ref{lem:energy-ineq} with $\alpha_1=\AlphaOne$, $u=0.0068682$ and $c_2=\Ctwo$. They also give the
inequality of the constrained case with the same $U_0$ and $c_2$.

\emph{The squares that remain.} Section~\ref{sec:signed} uses what the completion of squares leaves, so we record it for the case with doubled
interaction. Let $\mu_F$, $\mu_M$ and $\mu_S$ be the normalized counting measures of the free points, the prescribed points of the first kind and the
prescribed points of the second kind. Their total masses are $1-\alpha$, $\alpha-\alpha_2$ and $\alpha_2$. Put $W=\mu_F+\mu_M$,
$A_k(\mu)=\int T_k\,d\mu$ and $\langle a,b\rangle=\sum_ka_kb_k/k$, so that $\|a\|^2=\langle a,a\rangle$. Three products carry the interactions
in the variable $x$. The Vandermonde product of the free points comes from the Cauchy determinant. The row determinant gives the Vandermonde product of the
free points and the prescribed points of the first kind, and the Vandermonde product of the points $y$ of the second kind. These last points
enter the row determinant through $1/y$, so it also gives the cross factors $|w-1/y|=(1-wy)/y$, for $w$ a point of $W$.
Since $-\log|x-x'|=\log2+2\sum_kT_k(x)T_k(x')/k$ and $\log(1-wy)\le-wy$, the part of $E$ quadratic in the Chebyshev moments is, apart from the
diagonal terms and the multiples of $\log2$, at most
\[
 -\|A(\mu_F)\|^2-\|A(W)\|^2-\|A(\mu_S)\|^2-A_1(W)A_1(\mu_S).
\]
For any measure $\mu$ we have $-\|A(\mu)\|^2=\|u\|^2-2\langle u,A(\mu)\rangle-\|A(\mu)-u\|^2$. Using this for $\mu_F$ and $W$, and writing
$A_1(W)=u_1+(A_1(W)-u_1)$, we find that the display equals
\[
 2\|u\|^2-2\langle u,A(2\mu_F+\mu_M)\rangle-u_1A_1(\mu_S)-\mathcal Q,
\]
where
\[
 \mathcal Q=\|A(\mu_F)-u\|^2+\|A(W)-u\|^2+\|A(\mu_S)\|^2+(A_1(W)-u_1)A_1(\mu_S).
\]
The first three terms are the term $2\|u\|^2$ of $\mathcal N$ and the terms of $X$, $F_1$ and $F_2$ that contain $u$. So the energy bound obtained from
Lemma~\ref{lem:energy-ineq} improves by $\mathcal Q$. To bound $\mathcal Q$ from below, let $e_1$ be the first
unit vector and put $a=A(\mu_F)-u$ and $b=A(W)-u+\frac12A_1(\mu_S)e_1$. Since $\|e_1\|=1$,
\[
 \|b\|^2=\|A(W)-u\|^2+(A_1(W)-u_1)A_1(\mu_S)+\tfrac14A_1(\mu_S)^2 .
\]
Also $\frac14A_1(\mu_S)^2\le\|A(\mu_S)\|^2$, so $\mathcal Q\ge\|a\|^2+\|b\|^2\ge2\|\frac12(a+b)\|^2$. The average of $a$ and $b$ is
$A(\mu_F+\frac12\mu_M)-u+\frac14A_1(\mu_S)e_1$. This gives
\begin{equation}\label{eq:squares}
 \mathcal Q\ge\Delta,\qquad\Delta=2\sum_k\frac1k\bigl(A_k(\mu_F+\tfrac12\mu_M)+\tfrac14A_1(\mu_S)\delta_{k1}-u_k\bigr)^2 .
\end{equation}
The weight $\frac12$ of $\mu_M$ records that these points occur in one of the two Vandermonde products of the $x$ variables, and the weight
$\frac14$ is half of the $\frac12$ in $b$. The diagonal terms are treated by OpenAI's regularization, which replaces the moments by damped moments
with an error that tends to zero, as in the proof of Lemma~\ref{lem:freetangent}. We suppress this.

\subsection{The trial sequences}\label{ssec:trials}
The trial sequences have the form
\begin{equation}\label{eq:trial}
 u_k=c_k+\operatorname{Re}\sum_jr_jz_j^k\qquad(k\ge1;\ c_k=0\text{ for }k>\Kcoef),
\end{equation}
with $\Kcoef$ Chebyshev coefficients and finitely many geometric terms. The geometric terms have a simple
meaning. Since $\sum_{k\ge1}z^kT_k(x)/k=-\frac12\log(1-2xz+z^2)$, a term $(z,r)$ contributes
$-\frac12\operatorname{Re}\,r\log(1-2xz+z^2)$ to the potential $\mathsf T(u,x)$. If
$x_0+i\epsilon_0=\frac12(z+z^{-1})$, then $1-2xz+z^2=2z(x_0+i\epsilon_0-x)$. So the term is, up to a constant,
the logarithmic potential of a point charge of complex strength $r$ at $x_0+i\epsilon_0$ and its conjugate. The
extremal measure has a gap in its support around $x=0$. Its potential has a logarithmic singularity at $x=0$
and square-root behaviour at the edges of the gap. A Chebyshev expansion converges slowly near such points,
while a charge at a small distance $\epsilon_0$ reproduces the local behaviour with one term. This is also why
the quadrature nodes of Section~\ref{sec:real} avoid a neighbourhood of $x=0$. Appendix~\ref{app:field} gives the charges
and the certified bounds.

\subsection{The two cases at the improved rate}\label{ssec:cases}
Theorem~\ref{thm:real} asks for the rate $U_0=U+\Gamma_0/n_0^2$ instead of $U$. In the constrained case the
inequality above holds at $U_0$ for every $\alpha$, with the slope $c_2$. In the case with doubled interaction
Lemma~\ref{lem:energy-ineq} gives it at $U_0$ only for $\alpha\ge\alpha_1$, because the extra term $U_0-U$ is
absorbed by the term $c_2\alpha$ and the quadratic term once $\alpha$ is not small. So a new argument is needed
only in the case with doubled interaction with fewer than $\alpha_1n$ prescribed columns.
Section~\ref{sec:signed} supplies it, with the same constant $c_2$. Precisely, it proves the following.

\begin{proposition}\label{prop:signed-goal}
In the case with doubled interaction, for every configuration with $m\le\alpha_1n$ prescribed columns, the
integrand of Lemma~\ref{lem:andreief}, with the generalized Vandermonde determinant replaced by the bound of
Lemma~\ref{lem:binomvand}, is at most
\[
 \exp\Bigl(n^2\bigl(-U_0+\tfrac12\log2+c_2\alpha+o(1)\bigr)\Bigr)\,C_0^m\prod_j\mathcal B(k_j),\qquad \alpha=m/n,
\]
with an absolute constant $C_0$.
\end{proposition}

\begin{proof}[Proof of Theorem~\ref{thm:real}, given Proposition~\ref{prop:signed-goal}]
By Proposition~\ref{prop:signed-goal} and the discussion above, in both cases a determinant $M_f$ with $m$
prescribed columns $k_1,\ldots,k_m$ is at most
$C_0^m\prod_j\mathcal B(k_j)e^{n_0c_2Nm}\exp\bigl(n_0^2N^2(-U_0+\frac12\log2+o(1))\bigr)$. The factor $C_0^m$ comes
from the expansion of the row functions at the prescribed points. By \eqref{eq:wk} and the bound for
$\mathcal B(k)$, the constants of Lemma~\ref{lem:expansion} satisfy
\[
 \log(B_kC_k)\le\bigl(n_0c_2+w(k/n)+n_0\tilde h_2(k/n)+o(1)\bigr)N\le(\rho+o(1))N .
\]
For $|\varepsilon|\le e^{-\rho N}$ Lemma~\ref{lem:expansion} gives a factor $(1+O(N^5))^n=e^{o(N^2)}$. With
$D=n_0^2(U_0-\frac12\log2)$ this is the claim.
\end{proof}

\section{Proof of Theorem~\ref{thm:main}}\label{sec:proof}

\subsection{The criterion}
We prove Proposition~\ref{prop:criterion}. The proof shows where each of the three properties of $Q_N$
enters. The nonvanishing and the smallness are used once, at the end. The divisibility by squares of primes
is what makes the choice of the prime free.

\begin{proof}[Proof of Proposition~\ref{prop:criterion}]
Let $a/b$ be a rational number with $b$ large. Put $L=\log b$, $U=\kappa L$ and $V=(\kappa-\omega)L$, and
suppose $|\xi-a/b|\le b^{-\rho\kappa}$. By the prime number theorem, $(1-\delta)t\le\theta(t)\le(1+\delta)t$
for all large $t$, and $V\to\infty$ with $b$. If every prime $p\in(V,U]$ had $p^2\mid b$, then
$2(\theta(U)-\theta(V))\le\log b=L$. But
\[
 2\bigl(\theta(U)-\theta(V)\bigr)\ge2L\bigl((1-\delta)\kappa-(1+\delta)(\kappa-\omega)\bigr)
 =2L\bigl(\omega-\delta(2\kappa-\omega)\bigr)>L
\]
by the fourth hypothesis. So there are primes in $(V,U]$ whose squares do not divide $b$. Let $p$ be the
\emph{largest} one. Every prime $q\in(p,U]$ has $q^2\mid b$, and $q\le U<2V<2p$ because $2\omega<\kappa$.
Let $B$ be the product of these $q^2$. Then $B\mid b$ and $\log B=2(\theta(U)-\theta(p))$.

Consider the integer $b^{dp}Q_p(a/b)=\sum_ic_ia^ib^{dp-i}$. Each of its terms is divisible by $q^{2dp}$ for
each of these $q$. Indeed $q^{2i}$ divides $c_i$ by hypothesis, and $q^2\mid b$ gives $q^{2(dp-i)}\mid b^{dp-i}$.
So $B^{dp}$ divides the integer. Since $p^2\nmid b$, $Q_p(a/b)\ne0$. So $b^{dp}Q_p(a/b)/B^{dp}$ is a nonzero
integer, and
\begin{equation}\label{eq:recycled}
 1\le\Bigl(\frac bB\Bigr)^{dp}\,|Q_p(a/b)| .
\end{equation}
Since $\rho p\le\rho U=\rho\kappa L$, we have $|a/b-\xi|\le e^{-\rho p}$. So $|Q_p(a/b)|\le e^{-\eta p^2}$, and
\eqref{eq:recycled} gives $\eta p^2\le dp\bigl(L-2(\theta(U)-\theta(p))\bigr)$. With
$\theta(U)-\theta(p)\ge U-p-2\delta U$ this is
\[
 \eta p\le d\bigl(L-2U+2p+4\delta U\bigr).
\]
Hence $(\eta-2d)p\le dL-2dU+4d\delta U$. As $\eta<2d$ and $p\le U$, the left side is at least $(\eta-2d)U$. So
$\eta U\le dL+4d\delta U$, that is, $\eta\kappa\le d+4d\delta\kappa$. This contradicts the last hypothesis.
\end{proof}

The largest admissible prime is the right choice. The primes of the window above it have their squares in
$b$, and the divisibility hypothesis recovers these squares from the coefficients. A cruder choice, of any
prime in a short interval that does not divide $b$, loses a factor. The free lower end $V=(\kappa-\omega)L$
allows $\kappa<2$, because the condition $q\le U<2p$ for the recycled primes needs only $U<2V$.

\subsection{The hypotheses for the design $D$}
Take the design $D=(58,9,3,5,3)$ and $d=53$. The polynomials $Q_p$ of the proposition are the $Q_N$ with
$N=p$. Let
\begin{gather*}
 \rho=\Rho,\quad \eta=\EtaTot,\quad \omega=\tfrac35,\quad \delta=2.9\cdot10^{-7},\\
 \kappa=\frac{410}{\rho}=\Kappa\ldots
\end{gather*}
We check the hypotheses.
\begin{itemize}
\item $\deg Q_N\le53N$ by \eqref{eq:QN} and Lemma~\ref{lem:degree}.
\item The nonvanishing. By Proposition~\ref{prop:nonvanishing}, $\Delta_p(a/b)\ne0$ for large $p$ with
$p^2\nmid b$. Hence $Q_p(a/b)\ne0$.
\item The divisibility. By Lemma~\ref{lem:recycling}, $q^{2i}$ divides the $i$-th coefficient of $Q_p$ for all
primes $p<q\le2p$.
\item The smallness. By Theorem~\ref{thm:real}, Proposition~\ref{prop:denominators} and
$\log(53N)!=o(N^2)$, for $|x-G|\le e^{-\rho N}$ and large $N$,
\[
 |Q_N(x)|\le\exp\bigl(-(\TDelta\ldots-\Aden\ldots-3\varepsilon)N^2\bigr)\le e^{-\eta N^2},
\]
since $\eta=\EtaTot<\EtaExact\ldots$ In the formal proof this is the inequality
$\eta<n_0^2(U_0-\frac12\log2)-\mathsf A$ with $\int\kappa$ replaced by its lower bound $\ChiLean$.
\item The numerical conditions. $\rho\kappa=410$ and $0<\eta<2d=106$. Next $\omega=3/5\le1$ and
$2\omega=1.2<\kappa$. Further $2(\omega-\delta(2\kappa-\omega))=1.199998\ldots>1$. Finally
$d+4d\delta\kappa=\PayL\ldots<\eta\kappa=\PayR\ldots$
\end{itemize}
Proposition~\ref{prop:criterion} therefore gives $|G-a/b|>b^{-\rho\kappa}=b^{-410}$ for all integers $a$ and all
sufficiently large $b$. This proves Theorem~\ref{thm:main}. The last condition is equivalent to
$410>x=d\rho/(\eta-4d\delta)=\Bx\ldots$

\section{Remarks on the method}\label{sec:remarks}

\subsection{The three ingredients}
The exponent $d\rho/\eta$ has three ingredients. The degree $d=n-g$ comes from the construction. Of
$\rho=\Rho$, the energy cost of the fixed columns $n_0c_2$ is $208.62$, and the weight and the binomial factor
of one column contribute $\Psi_0=\PsiZero$. Bounding the two separately, by their maxima $52.79$ and
$58\log2=40.20$, would give $\rho=\RhoOld$. The decay rate $\eta$ is a small difference of two large numbers.
The rank deficiencies of Section~\ref{sec:content} contribute $22.05$ of $\eta=\EtaTot$, and the signed bound
of Section~\ref{sec:signed} contributes $\Gain$. Since the exponent is inversely proportional to $\eta$, a further
gain of $0.1$ in $\eta$ lowers it by about $1.1$.

\subsection{The determinant of the row functions}
The energy method bounds $|\det\mathsf L|$ by $6^n$, which is $e^{O(N)}$ and invisible at the scale $N^2$. At
the extremal configuration, numerical evaluation for small $N$ suggests that $\log|\det\mathsf L|$ is about
$-28N^2$, so most of the true size is not captured. The bound \eqref{eq:eleven} loses only $e^{O(N)}$, and
Theorem~\ref{thm:heine} is an identity. The loss is in the log-gas bounds of \S\ref{ssec:loggas}, which use fixed
trial measures, and in the requirement that the combined functional be bounded by its maximum over all
configurations. A sharper treatment of the denominator $Z_s(\psi)$ would raise the gain. These numerical values
are not part of the formal proof.

\subsection{Designs}
All the results of Section~\ref{sec:content} are lower bounds for the nullity, and the proofs give the kernel
vectors explicitly. The closed formulas, the counting inequality and the families derived from it are proved
for every design. So are the other arguments of the paper, except the stored subspaces of
Section~\ref{ssec:packets} and the trials, which belong to the design $D$. A different design changes $d$,
$\rho$ and $\eta$ together. We make no claim that $D$ is optimal. Except for one pair of families, the function
$\kappa$ uses no relation between different families. At a prime where windows overlap the exact nullity may
be larger than each bound. An upper bound for the nullity is not needed. More independence between families,
or larger stored subspaces, would increase $\eta$ without changing anything else.

\section{Formalization}\label{sec:formal}

All results of this paper are formalized in Lean~4 \cite{Lean4} with Mathlib \cite{Mathlib}, on top of OpenAI's
formalization of the irrationality of $G$ \cite{OpenAI2026}. The development is available at
\begin{center}
\FormalizationLink.
\end{center}
Theorem~\ref{thm:main} is the theorem
\begin{center}
{\small\texttt{CatalanMeasure.GenD.Signed.catalan\_irrationalityMeasure\_410}}
\end{center}
in \lean{CatalanMeasure/GenD/Signed/Data/Assembly/CellsAll.lean}. Its statement uses only Mathlib. Catalan's
constant is written as the series $\sum_j(-1)^j/(2j+1)^2$, and the conclusion is $|G-a/b|>b^{-410}$ for all
integers $a$ and all sufficiently large $b$. The modules contain no \lean{sorry} and no \lean{native\_decide}.
The command \lean{\#print axioms} checks that only the three standard axioms \lean{propext},
\lean{Classical.choice} and \lean{Quot.sound} are used. The repository also contains a configuration
\lean{comparator.json} for \lean{lake comparator} \cite{Comparator}. It checks that a solution file proves a
challenge statement with these axioms only. The prime number theorem is taken from the PNT+ project
\cite{PNTPlus}. The parts of OpenAI's formalization that are used were converted to Lean's module system and
ported to a recent Lean release. The repository's \texttt{README} and \texttt{PORTING.md} describe the changes.

\subsection{Module layout}
The design-generic part of the development is in the directory \lean{GenD} of \lean{CatalanMeasure}. A design is
data. The generic modules prove every statement for an arbitrary design, and the instance files contain
numerals only.
\begin{description}[leftmargin=1.5em,style=nextline]\sloppy
\item[Generic layer] \lean{PL} (piecewise linear profiles and their areas), \lean{PrimeGrowth} (the weighted
prime number theorem \eqref{eq:weighted-pnt}), \lean{Budget} (the scalar budgets), \lean{PayW} and
\lean{Recycled} (Proposition~\ref{prop:criterion}), and \lean{Final} (the structure \lean{DesignData} and the
theorem \lean{measure\_of\_design}, which contain no numeral).
\item[Denominators] \lean{OddProfile}, \lean{TwoAdic}, \lean{DenomExp}, \lean{DenomGrowth},
\lean{Denominator} and \lean{Normalizer}, with \lean{Uniform/Arithmetic} and \lean{Uniform/LedgerW} (the
common denominator built from an arbitrary lower bound for the nullity, and the growth of its logarithm).
\item[Rank deficiencies] \lean{RankT}, \lean{RankU}, \lean{RankU6}, \lean{RankS} and their shared layers for
Theorem~\ref{thm:content}; \lean{Uniform/W/} for the counting inequality (\lean{Core}, \lean{Profile}),
Lucas's step (\lean{Lucas/PStep}), the families of Section~\ref{ssec:families} and the stored data of
Section~\ref{ssec:packets} (\lean{Prof/Packet}, \lean{Charts6/Check}); the profile of $D$ with its area
\lean{gain532} and the bound \lean{gain532\_ge} in \lean{Uniform/Bank532}.
\item[Nonvanishing] \lean{Nonvanishing}, with the residue-table certificates modulo $1009$.
\item[Real place, perturbation] \lean{Analytic/} (the energy method with fixed points, the column expansion
and the quadrature nodes); \lean{BinomVandermonde} (Lemma~\ref{lem:binomvand}) and \lean{Keyhole}
(Lemma~\ref{lem:keyhole}) in \lean{GenD/Signed}; \lean{D58/Budget590} (the budgets at the rates $U$ and $U_0$,
and $\rho=$ \lean{rhoEnd}); \lean{Signed/Data/S1} (the $14$ cells of the bound for $\Psi_0$).
\item[Real place, signed bound] \lean{GenD/Signed}: \lean{PickDet} and \lean{RealPlace}
(Lemma~\ref{lem:pick}, Theorem~\ref{thm:C3}, \eqref{eq:eleven}), \lean{LemmaT}, \lean{LemmaTFree} and
\lean{HalfDisc} (Lemma~\ref{lem:T}), \lean{RatioIdentity} and \lean{BernsteinSzego}
(Theorem~\ref{thm:heine}), \lean{TiltedTangent} and \lean{ArcsineDenominator} (Lemmas~\ref{lem:LGup}
and~\ref{lem:LGlow}), \lean{Spike} and its companions (Theorem~\ref{thm:spike}), \lean{ReflectionEnergy}
and \lean{ArcsinePhantom} (Theorem~\ref{thm:reflection}), \lean{MarkedExpansion}, \lean{SigmaCap} and
\lean{EnergyNear} (\S\ref{ssec:fixedrows}), \lean{LemmaK} and \lean{Brackets} (\S\ref{ssec:class}),
\lean{Interp} and \lean{Phase2Contract} (the interpolation of trials), and \lean{Gate} and \lean{Theorem410}
(the assembly).
\item[Trials and cells] \lean{Signed/Data/Trial} (the three trials), \lean{Signed/Data/Consts} (the constants
of the trials, evaluated in exact arithmetic), \lean{Signed/Data/Cells} (the cell engine and the cells of the
three trials) and \lean{Signed/Data/Assembly} (the records of the cell statements and the final theorem).
\end{description}

\subsection{Certificates}
The proof combines conventional formal proofs with numerical certificates checked by the Lean kernel. They
are the inverses of the base matrices modulo $1009$, the trial sequences and covers of
Sections~\ref{sec:real} and~\ref{sec:signed}, the cells of the bound for $\Psi_0$, and the stored subspaces,
kernel vectors and window conditions of Section~\ref{ssec:packets}. All are exact rational or integer
computations, decided by the kernel without compiled code. The interval arithmetic used to generate the trial
data and the covers \cite{Arb} is not trusted. Lean rechecks every inequality. Five devices keep the
certificates small enough for the kernel.
\begin{itemize}
\item \emph{Shared Taylor bounds.} The cells of a cover are grouped in blocks. For each block the derivatives of
the polynomial part of the field are evaluated once, exactly, at the centre of the block, together with one
bound for the remainder. Every cell of the block obtains its expansion by an exact shift.
\item \emph{Split shifts.} In the covers of Section~\ref{sec:signed}, only the low-order coefficients of the
shifted expansion are computed exactly. The others are bounded by one crude table per block.
\item \emph{Accuracy per term.} The logarithms are evaluated by the series
$\log\frac{1+x}{1-x}=2x\sum_ix^{2i}/(2i+1)$ with an explicit geometric tail. The number of terms is chosen for
each logarithmic term separately, from a table of nodes shared by all the cells.
\item \emph{Residue recurrences.} The tables of moments modulo $1009$ are generated by the one-step recurrences
that define them, once proved correct. The kernel checks the recurrence and the products of the matrices with
their claimed inverses, row by row, in natural numbers.
\item \emph{$2$-adic normalisation.} Where a recurrence or a shift produces integers divisible by large powers
of $2$, the tables store the quotients by these powers, with one lemma relating the two. This also avoids
integers that the kernel's hashing of literals treats as equal.
\end{itemize}

\section*{Disclosure of AI use}

The mathematics, the Lean formalization, the certificate generators and the text of this paper were produced by
AI systems at the request of the author, who posed the problems, directed the work and is responsible for this
paper: Claude Opus~5.5, Claude Sonnet~5.5 and Claude Haiku~5.5 (Anthropic), running in Claude Code, and
GPT-6.1 Sol (OpenAI), running in Codex. The underlying irrationality proof for $G$ and its Lean formalization are
OpenAI's \cite{OpenAI2026}.

\subsection*{Acknowledgements}
We thank the authors of Mathlib and of the PNT+ project, and OpenAI for making its formalization available.


\bibliographystyle{amsplain}
\bibliography{refs}

\end{document}
